\subsubsection{Exponential invariant system: unification of Pontryagin negative squares, Toeplitz failure patterns, endpoint/basepoint tomography, and Tate/Maslov obstructions}\label{subsubsec:xi-pontryagin-toeplitz-blaschke-tomography-tate}
This section unifies four classes of "failure/residue" objects appearing in different frameworks as the same integer exponent:
\emph{off-axis zeros (Blaschke point defects)}, \emph{failure patterns at the Toeplitz--PSD hierarchy}, \emph{invertible tomography via endpoint/basepoint scanning}, and \emph{\(\ZZ/2\)-equivariant uniformization obstructions (Tate/Maslov indices)}.
The structural point is: under the auditable framework of no singular inner factor and finite defects, all off-critical-slice point information is carried by a finite Blaschke factor \(B_{F_{\zeta}}\); this factor simultaneously controls
\begin{enumerate}
  \item Stable negative squares (Pontryagin index) of Toeplitz quadratic forms,
  \item Point spectrum residual axis in the functional model (model space \(K_{B_{F_\zeta}}\) / compressed shift \(A_{B_{F_\zeta}}\)),
  \item Pole data from Poisson--Busemann flux at endpoint \(-1\) and basepoint scanning profiles,
  \item And the point-like (Maslov/Toeplitz) integer part of the equivariant Tate obstruction (Proposition \ref{prop:selfdual-maslov-singular-dichotomy}).
\end{enumerate}

\paragraph{Setup and notation}
Following the Hardy/Smirnov conditional framework in Appendix \S\ref{subsec:unit-disk-inner-outer} and the Blaschke defect notation in \S\ref{subsubsec:selfdual-blaschke-defect-fixed-slice}.
Assume the disk completion readout \(F_{\zeta}\) is of bounded type and has no singular inner factor (\(S_{F_{\zeta}}\equiv 1\)), and its Blaschke factor \(B_{F_{\zeta}}\) is a finite Blaschke product.
Denote
\[
\kappa:=\deg(B_{F_{\zeta}})=\mathfrak d_{\mathrm{Bl}}(F_{\zeta})=\dim K_{B_{F_{\zeta}}}
\qquad\text{(Definition \ref{def:blaschke-defect-dimension}, Remark \ref{rem:blaschke-defect-finite})}.
\]
On the other hand, the horizon Carathéodory logarithmic derivative function \(\mathscr{C}_{\zeta}\) has power series at \(w=0\) written as (Appendix Theorem \ref{thm:app-horizon-toeplitz-lmi})
\[
\mathscr{C}_{\zeta}(w)=c_0+2\sum_{n\ge 1}c_n w^n,
\qquad c_{-n}:=\overline{c_n},
\]
and define the Toeplitz principal submatrix
\(
\mathbf{T}_N:=(c_{j-k})_{0\le j,k\le N}\in\Mat_{N+1}(\CC)_{\mathrm{Herm}}
\).
For a Hermitian matrix \(M\), denote \(\nu_-(M)\) as the multiplicity of its strictly negative eigenvalues (negative inertia index).

\paragraph{Pontryagin indexing: Toeplitz stable negative squares coincide with Blaschke defect}
From \(\mathscr{C}_{\zeta}\) define the Hermitian kernel
\begin{equation}\label{eq:xi-pontryagin-carath-kernel}
\boxed{\;
K_{\mathscr{C}_{\zeta}}(w,\zeta)
:=\frac{\mathscr{C}_{\zeta}(w)+\overline{\mathscr{C}_{\zeta}(\zeta)}}{1-w\overline{\zeta}}
\qquad(w,\zeta\in\mathbb{D}).\;}
\end{equation}

\begin{theorem}[Toeplitz negative squares index \(=\) Blaschke defect dimension]\label{thm:xi-toeplitz-negative-squares-equals-blaschke-defect}
Under the above finite defect setting, the kernel \(K_{\mathscr{C}_{\zeta}}\) has \textbf{exactly \(\kappa\) negative squares} on \(\mathbb{D}\) (Pontryagin index is \(\kappa\)).
Equivalently, the negative inertia index of Toeplitz principal submatrices satisfies
\[
\nu_-(\mathbf{T}_N)\le \kappa\qquad(\forall\,N\ge 0),
\qquad
\nu_-(\mathbf{T}_N)=\kappa\qquad(N\ge N_0(\kappa)).
\]
In particular, RH is equivalent to \(\kappa=0\), i.e., \(K_{\mathscr{C}_{\zeta}}\) is positive definite, and \(\mathbf{T}_N\succeq 0\) for all \(N\) (Appendix Theorem \ref{thm:app-horizon-toeplitz-lmi}).
\end{theorem}

\begin{proof}[Proof sketch (structural reduction)]
Under RH, \(\mathscr{C}_{\zeta}\) is a Carathéodory function, \eqref{eq:xi-pontryagin-carath-kernel} is a positive definite kernel and corresponds to \(\mathbf{T}_N\succeq 0\) (Appendix Theorem \ref{thm:app-horizon-toeplitz-lmi}).
If there exist in-disk point defects, then \(F_{\zeta}\) has \(\kappa\) zeros in \(\mathbb{D}\) (counted with multiplicity), i.e., \(B_{F_{\zeta}}\) is a degree-\(\kappa\) finite Blaschke product;
since \(\mathscr{C}_{\zeta}(w)=-\Xi'_{\zeta}(s(w))/\Xi_{\zeta}(s(w))\), these zeros induce \(\kappa\) in-disk poles on \(\mathscr{C}_{\zeta}\), thus \(K_{\mathscr{C}_{\zeta}}\) is a positive definite kernel minus a rank-\(\kappa\) modification, generating exactly \(\kappa\) negative squares.
Taking the Gram matrix of this kernel on the finite-dimensional subspace spanned by \(\{1,w,\dots,w^N\}\) gives a family of Hermitian matrices congruent to \(\mathbf{T}_N\); congruence does not change inertia index, thus yielding the stated stable negative squares conclusion.
(This "finite poles \(\Leftrightarrow\) finite negative squares" phenomenon is compatible with the finite-order refuter language in \S\ref{subsubsec:xi-verblunsky-cmv-from-local-jet}.)
\end{proof}

\paragraph{Canonical generation of negative directions: model space gives Toeplitz failure witnesses}
The finite Blaschke factor not only gives the "number of negative squares," but also provides the \emph{source} of negative directions: they are forced by the point spectrum modes of the model space \(K_{B_{F_\zeta}}\) (Theorem \ref{thm:blaschke-point-spectrum-correspondence}).

\begin{theorem}[Minimal generation of failure witnesses: negative subspace isomorphic to point spectrum residual axis]\label{thm:xi-toeplitz-witness-generated-by-model-space}
Under the setting of Theorem \ref{thm:xi-toeplitz-negative-squares-equals-blaschke-defect}, there exists a family of linear embeddings stabilizing with \(N\)
\[
\iota_N:K_{B_{F_{\zeta}}}\hookrightarrow \CC^{N+1},
\]
such that the Toeplitz quadratic form \(x\mapsto x^\ast \mathbf{T}_N x\) restricted to the subspace \(\iota_N(K_{B_{F_{\zeta}}})\) has \textbf{exactly \(\kappa\)-dimensional negative definite part} (when \(N\ge N_0\)).
Moreover, when \(N\) is sufficiently large, this negative definite part can be generated by the point modes of the point spectrum generating subspace
\(\mathcal{H}_{\mathrm{pp}}(A_{B_{F_{\zeta}}})\)
(Theorem \ref{thm:blaschke-point-spectrum-correspondence}) via the image of \(\iota_N\).
\par
Therefore, all negative directions of the Toeplitz hierarchy are not diffusely generated, but are forced by \(B_{F_{\zeta}}\); its minimal dimension is exactly \(\deg(B_{F_{\zeta}})=\kappa\), which is the "residual axis minimal orthogonal dimension" in the regime semantics.
\end{theorem}

\begin{proof}[Proof sketch]
By Theorem \ref{thm:xi-toeplitz-negative-squares-equals-blaschke-defect}, the negative inertia index of \(\mathbf{T}_N\) stabilizes at \(\kappa\).
On the other hand, \(\kappa=\dim K_{B_{F_\zeta}}=\dim\mathcal H_{\mathrm{pp}}(A_{B_{F_\zeta}})\) (Theorem \ref{thm:blaschke-point-spectrum-correspondence}).
Take \(\iota_N\) as the truncated coordinate map of \(K_{B_{F_\zeta}}\) on the Hardy basis \(\{1,w,w^2,\dots\}\) (e.g., taking the first \(N{+}1\) Taylor coefficients), which gives a sequence of explicitly auditable embeddings; under these coordinates, the Toeplitz quadratic form coincides with the kernel Gram form, thus its negative subspace can be identified equivalently as the negative direction of the model space image.
Point mode generation is given by \(\mathcal H_{\mathrm{pp}}(A_B)=K_B\) in the finite-dimensional case and \(\deg B=\dim K_B\). Proven.
\end{proof}

\paragraph{Operatorization of endpoint coupling strength: Poisson--Busemann weight equals sum of squares of point mode boundary couplings}
The Poisson kernel weight at endpoint \(-1\), \(P_a(-1)=\frac{1-|a|^2}{|1+a|^2}=\mathcal{B}_{-1}(a)\) (\eqref{eq:app-poisson-kernel-endpoint-minus1}), is not only a geometric characterization but can also be interpreted as the coupling strength of point modes to endpoint readout.

\begin{theorem}[Operatorization of endpoint flux: Julia index, trace identity, and point mode coupling sum of squares]\label{thm:xi-endpoint-flux-pointmode-coupling}
Let \(B\) be a finite Blaschke product with zero multiset \(\{a_k\}_{k=1}^{\kappa}\subset\mathbb{D}\).
Denote the endpoint flux (absorption coefficient)
\[
\Phi_{-1}(B):=\sum_{k=1}^{\kappa}P_{a_k}(-1)
\qquad\text{(Proposition \ref{prop:xi-endpoint-absorption-coefficient})}.
\]
\par
Let the model space \(K_B:=H^2\ominus BH^2\), and define the endpoint evaluation functional \(L_{-1}:K_B\to\CC\) as \(L_{-1}(f):=f(-1)\).
By Theorem \ref{thm:blaschke-point-spectrum-correspondence}, in the finite Blaschke case, \(\mathcal{H}_{\mathrm{pp}}(A_B)=K_B\), thus \(L_{-1}\) can also be viewed as endpoint evaluation on the point spectrum residual axis.
Since \(K_B\) is finite-dimensional, \(L_{-1}\) is bounded, and its Julia--Carathéodory index
\(
\mathcal{J}_{-1}(B):=\lim_{r\uparrow 1}\frac{1-|B(-r)|^2}{1-r^2}
\)
satisfies (Theorem \ref{thm:xi-endpoint-julia-indicator-equals-absorption})
\[
\boxed{\;
\Phi_{-1}(B)=\mathcal{J}_{-1}(B)=\mathrm{tr}\!\bigl(L_{-1}^\ast L_{-1}\bigr)\; }.
\]
On the other hand, for the Hardy reproducing kernel vector \(k_{a}(w):=(1-\overline a\,w)^{-1}\) (which is also a point mode reproducing kernel of model space \(K_B\) when \(a\) is a zero of \(B\)), there holds the strict identity
\[
\boxed{\;
\Phi_{-1}(B)
=\sum_{k=1}^{\kappa}\frac{|k_{a_k}(-1)|^2}{\|k_{a_k}\|_{H^2}^{2}}
=\sum_{k=1}^{\kappa}\frac{|1+\overline{a_k}|^{-2}}{(1-|a_k|^2)^{-1}}
=\sum_{k=1}^{\kappa}\frac{1-|a_k|^2}{|1+a_k|^2}\; }.
\]
Therefore, \(\Phi_{-1}(B)\) equals the sum of squares of coupling strengths from each point mode on the point spectrum residual axis to endpoint readout, thus strictly formalizing "projection of off-axis information to endpoint" as a finite-sum trace-type invariant.
\end{theorem}

\begin{proof}
Since \(K_B\) is a finite-dimensional reproducing kernel space, the Riesz representation of endpoint evaluation \(L_{-1}\) coincides with the Julia index: for the model space reproducing kernel at \(\zeta\in\mathbb{D}\)
\[
k_\zeta^B(w):=\frac{1-\overline{B(\zeta)}B(w)}{1-\overline{\zeta}w}\in K_B
\]
satisfying
\(
\|k_\zeta^B\|_{K_B}^2=k_\zeta^B(\zeta)=\frac{1-|B(\zeta)|^2}{1-|\zeta|^2}
\).
Taking \(\zeta=-r\uparrow-1\) and using radial continuability of finite Blaschke, we obtain
\(
\|k_{-1}^B\|_{K_B}^2=\mathcal{J}_{-1}(B)
\);
and \(L_{-1}^\ast L_{-1}\) is a rank-one operator whose trace satisfies
\(
\mathrm{tr}(L_{-1}^\ast L_{-1})=\|k_{-1}^B\|_{K_B}^2
\).
Again by Theorem \ref{thm:xi-endpoint-julia-indicator-equals-absorption}, \(\mathcal{J}_{-1}(B)=\Phi_{-1}(B)\).
\par
For the point mode coupling sum of squares identity, directly compute
\(
k_{a}(-1)=\frac{1}{1+\overline a}
\), and \(\|k_a\|_{H^2}^2=\frac{1}{1-|a|^2}\).
Substituting and summing over \(a=a_k\) yields the conclusion. Proven.
\end{proof}

\paragraph{Full-matrix limit of negative spectrum: Pick--Poisson Gram matrix and pseudohyperbolic Vandermonde invariant}
The previous paragraph interprets endpoint flux \(\Phi_{-1}(B)\) as the sum of squares of point mode coupling strengths to endpoint readout (Theorem \ref{thm:xi-endpoint-flux-pointmode-coupling}), thus only reading out the \emph{diagonal} total quantity (trace invariant) of the negative spectrum.
This paragraph further indicates: in the semiclassical limit, the entire negative spectrum of the Toeplitz certificate is not encoded only by diagonal Poisson weights but is fully encoded by an explicit Pick--Poisson Gram matrix; its determinant gives a complete invariant that is simultaneously sensitive to "near-boundary/collision."

\begin{proposition}[Full-matrix limit of negative spectrum matrix: lifting from diagonal Poisson weights to Pick--Poisson Gram matrix]\label{prop:xi-toeplitz-negative-spectrum-pick-poisson-gram}
Let the point defects be given by the finite Blaschke zero multiset \(\{a_j\}_{j=1}^{\kappa}\subset\mathbb{D}\) (expanded by multiplicity), and let \(\mathbf{T}_N\) be the corresponding Carathéodory/Toeplitz principal submatrix (\(\mathbf{T}_N=(c_{j-k})_{0\le j,k\le N}\)).
For each \(a\in\mathbb{D}\), define the truncated reproducing vector
\[
v_a^{(N)}:=(1,\overline a,\overline a^{\,2},\dots,\overline a^{\,N})^{\mathsf T}\in\CC^{N+1},
\qquad
\mathcal{K}_N:=\Span\{v_{a_j}^{(N)}:1\le j\le\kappa\}\subset\CC^{N+1}.
\]
Then there exists a family of linear isomorphisms \(R_N:\CC^\kappa\to\mathcal{K}_N\) (induced by fixed basis choice) such that as \(N\to\infty\) and \(N(1-|a_j|)\to\infty\) (for all \(j\)), there holds operator norm convergence
\[
\boxed{\;
R_N^\ast\bigl(\mathbf{T}_N|_{\mathcal{K}_N}\bigr)R_N\ \longrightarrow\ -\,\mathsf{P}\;}
\]
where the limit matrix \(\mathsf{P}\in\Mat_{\kappa}(\CC)_{\mathrm{Herm}}\) is the Pick--Poisson Gram matrix
\begin{equation}\label{eq:xi-pick-poisson-gram}
\boxed{\;
\mathsf{P}_{jk}
:=
\frac{(1-|a_j|^2)(1-|a_k|^2)}{(1+\overline a_j)(1+a_k)(1-\overline a_j a_k)}
\qquad(1\le j,k\le\kappa).\;}
\end{equation}
This matrix is positive semi-definite; when \(\{a_j\}\) are pairwise distinct, it is positive definite.
Moreover, diagonal entries satisfy
\[
\mathsf{P}_{jj}=\frac{1-|a_j|^2}{|1+a_j|^2}=P_{a_j}(-1),
\]
thus \(\mathrm{tr}(\mathsf{P})=\sum_{j=1}^{\kappa}P_{a_j}(-1)=\Phi_{-1}(B)\) consistent with the trace identity in Theorem \ref{thm:xi-endpoint-flux-pointmode-coupling}.
\end{proposition}
\begin{proof}[Proof sketch]
By \(\langle k_{a_j},k_{a_k}\rangle_{H^2}=(1-\overline a_j a_k)^{-1}\) and the weight factor
\(
u(a):=\frac{1-|a|^2}{1+\overline a}
\),
one can write \(\mathsf P\) as the Gram of weighted reproducing kernels:
\(
\mathsf P_{jk}=\langle u(a_j)k_{a_j},u(a_k)k_{a_k}\rangle_{H^2}
\),
thus positive semi-definiteness is immediate.
On the other hand, under the finite defect framework, the Toeplitz negative subspace is generated by model space point modes in Hardy basis truncation (Theorem \ref{thm:xi-toeplitz-witness-generated-by-model-space}), and when \(N(1-|a_j|)\to\infty\), truncation error in operator norm sense can be controlled by \(|a_j|^{N}\) and vanishes; restricting the Toeplitz quadratic form to \(\mathcal K_N\) and normalizing under a fixed basis yields the stated limit. Proven.
\end{proof}

\begin{proposition}[Complete invariant via negative spectrum determinant: factorization law of Poisson weights and pseudohyperbolic Vandermonde]\label{prop:xi-pick-poisson-det-factorization}
Following the setting of Proposition \ref{prop:xi-toeplitz-negative-spectrum-pick-poisson-gram}.
Let the pseudohyperbolic distance
\[
\rho(a,b):=\left|\frac{a-b}{1-\overline b a}\right|\qquad(a,b\in\mathbb{D}).
\]
Then the determinant of the Pick--Poisson Gram matrix \(\mathsf P\) has the closed-form factorization
\[
\boxed{\;
\det\mathsf P
=\Bigl(\prod_{j=1}^{\kappa}\frac{1-|a_j|^2}{|1+a_j|^2}\Bigr)
\Bigl(\prod_{1\le j<k\le\kappa}\rho(a_j,a_k)^2\Bigr).\;}
\]
Moreover, let \(\{\lambda_1^-(\mathbf{T}_N),\dots,\lambda_\kappa^-(\mathbf{T}_N)\}\) be the \(\kappa\) negative eigenvalues of \(\mathbf{T}_N\) (counted with multiplicity, existing for \(N\gg 1\) and guaranteed by Theorem \ref{thm:xi-toeplitz-negative-squares-equals-blaschke-defect}), then there holds the product limit
\[
\boxed{\;
\lim_{N\to\infty}\ \prod_{j=1}^{\kappa}\bigl(-\lambda_j^{-}(\mathbf{T}_N)\bigr)\ =\ \det\mathsf P. \;}
\]
This product law degenerates strictly for both zero collisions (\(\rho(a_j,a_k)\to 0\)) and near-boundary (\(|a_j|\uparrow 1\)), thus giving a global complete invariant of point defects (in the semiclassical limit sense).
\end{proposition}
\begin{proof}[Proof sketch]
Write \(\mathsf P\) as the congruence transformation of diagonal scaling and standard reproducing kernel Gram:
\[
\mathsf P
=D\,G\,D^\ast,
\qquad
D:=\mathrm{diag}\!\left(\frac{1-|a_j|^2}{1+\overline a_j}\right)_{j=1}^{\kappa},
\qquad
G_{jk}:=\frac{1}{1-\overline a_j a_k}.
\]
By the standard formula for Cauchy determinant/kernel Gram, \(\det G=\bigl(\prod_j(1-|a_j|^2)^{-1}\bigr)\prod_{j<k}\rho(a_j,a_k)^2\).
Combining with \(\det D=\prod_j\frac{1-|a_j|^2}{1+\overline a_j}\) yields the stated factorization.
The product limit follows from the operator norm convergence in Proposition \ref{prop:xi-toeplitz-negative-spectrum-pick-poisson-gram}: negative spectrum on \(\mathcal K_N\) converges continuously to \(-\mathsf P\), thus the product of negative eigenvalues converges to \(\det\mathsf P\). Proven.
\end{proof}

\begin{corollary}[Explicit lower bound on Pick--Poisson minimal spectral gap: sealed by endpoint flux and pseudohyperbolic separation]\label{cor:xi-pick-poisson-gap-lower-bound}
Following the setting of Propositions \ref{prop:xi-toeplitz-negative-spectrum-pick-poisson-gram}--\ref{prop:xi-pick-poisson-det-factorization}, denote
\[
p_j:=\mathsf P_{jj}=P_{a_j}(-1)=\frac{1-|a_j|^2}{|1+a_j|^2},
\qquad
\rho_{jk}:=\rho(a_j,a_k).
\]
Then \(\mathsf P\) is positive definite and its minimal eigenvalue satisfies the auditable lower bound
\[
\boxed{\;
\lambda_{\min}(\mathsf P)\ \ge\ \frac{\det\mathsf P}{\bigl(\mathrm{tr}\,\mathsf P\bigr)^{\kappa-1}}
\ =\
\frac{\Bigl(\prod_{j=1}^{\kappa}p_j\Bigr)\Bigl(\prod_{1\le j<k\le\kappa}\rho_{jk}^2\Bigr)}
{\Bigl(\sum_{j=1}^{\kappa}p_j\Bigr)^{\kappa-1}}\; }.
\]
Therefore, as long as endpoint flux does not collapse (\(\sum_j p_j>0\)) and spectrum points remain separated (\(\rho_{jk}>0\)), the spectral gap of \(\mathsf P\) has a strictly positive lower bound; when any \(p_j\to0\) or \(\rho_{jk}\to0\), this lower bound degenerates strictly, giving a quantitative certificate of "near-boundary/collision \(\Rightarrow\) spectral gap collapse."
\end{corollary}
\begin{proof}
Let \(\lambda_1\le\cdots\le\lambda_\kappa\) be the eigenvalues of \(\mathsf P\), then
\(
\det\mathsf P=\prod_{i=1}^{\kappa}\lambda_i
\le \lambda_1\,(\sum_{i=1}^{\kappa}\lambda_i)^{\kappa-1}
=\lambda_{\min}(\mathsf P)\,(\mathrm{tr}\,\mathsf P)^{\kappa-1}
\).
Rearranging gives the first formula; the second follows from the determinant factorization in Proposition \ref{prop:xi-pick-poisson-det-factorization} and \(\mathrm{tr}\,\mathsf P=\sum_j p_j\). Proven.
\end{proof}

\begin{proposition}[Negative spectrum margin for finite-order certificates: stable visibility threshold of Toeplitz negative eigenvalues]\label{prop:xi-toeplitz-negative-spectrum-margin-threshold}
Following the notation of Proposition \ref{prop:xi-toeplitz-negative-spectrum-pick-poisson-gram}, let
\[
A_N:=R_N^\ast\bigl(\mathbf{T}_N|_{\mathcal K_N}\bigr)R_N\in\Mat_{\kappa}(\CC)_{\mathrm{Herm}},
\qquad
\varepsilon_N:=\|A_N+\mathsf P\|_{\mathrm{op}}.
\]
If for some \(N\) there holds
\[
\varepsilon_N\ <\ \frac12\,\lambda_{\min}(\mathsf P),
\]
then the stable margin conclusion holds:
\[
\boxed{\;
\nu_-\!\bigl(\mathbf{T}_N\bigr)=\kappa,
\qquad
\lambda_i^{-}(\mathbf{T}_N)\ \le\ -\frac12\,\lambda_{\min}(\mathsf P)\quad(1\le i\le \kappa).\;}
\]
Therefore, finite-order Toeplitz certificates not only reveal the defect index \(\kappa\) in dimension but also possess a uniform negative margin in spectral amplitude; this margin is completely sealed by the algebraic lower bound of endpoint Poisson weights and pseudohyperbolic separation (Corollary \ref{cor:xi-pick-poisson-gap-lower-bound}).
\end{proposition}
\begin{proof}[Proof sketch]
By Weyl's inequality, eigenvalues of \(A_N\) differ from those of \(-\mathsf P\) by at most \(\varepsilon_N\) pointwise.
Therefore, when \(\varepsilon_N<\lambda_{\min}(\mathsf P)/2\), \(A_N\) is negative definite and all its eigenvalues are \(\le-\lambda_{\min}(\mathsf P)/2\).
This means the quadratic form \(x\mapsto x^\ast\mathbf{T}_N x\) on \(\mathcal K_N\) is strictly negative definite with uniform margin, thus \(\mathbf{T}_N\) has at least \(\kappa\) negative eigenvalues and its \(\kappa\)-th smallest eigenvalue does not exceed \(-\lambda_{\min}(\mathsf P)/2\).
On the other hand, under finite defect framework, the Toeplitz negative squares index is \(\kappa\), so for any \(N\) we must have \(\nu_-(\mathbf{T}_N)\le \kappa\); combining yields \(\nu_-(\mathbf{T}_N)=\kappa\).
Thus all \(\kappa\) negative eigenvalues of \(\mathbf{T}_N\) are at most \(-\lambda_{\min}(\mathsf P)/2\). Proven.
\end{proof}

\begin{corollary}[Closed-form threshold for visibility complexity: minimal order determined jointly by shallowest depth and weakest separation]\label{cor:xi-toeplitz-visibility-complexity-threshold}
Following the setting of Proposition \ref{prop:xi-toeplitz-negative-spectrum-margin-threshold}, denote the shallowest depth
\(
h_{\min}:=\min_j(1-|a_j|).
\)
If there exist constants \(c,C_0>0\) such that the error term satisfies
\[
\varepsilon_N\ \le\ C_0\,\exp(-cNh_{\min})\qquad(N\ge 1),
\]
then there exists a constant \(C>0\) such that when
\[
\boxed{\;
N\ \ge\ \frac{C}{h_{\min}}\Biggl[
\log(1+C_0)\ +\ (\kappa-1)\log\!\Bigl(\sum_{j=1}^{\kappa}p_j\Bigr)
\ -\ \sum_{j=1}^{\kappa}\log p_j
\ -\ 2\!\!\sum_{1\le j<k\le\kappa}\!\!\log\rho_{jk}
\Biggr]\;}
\]
it necessarily holds that \(\nu_-(\mathbf{T}_N)=\kappa\).
This threshold unifies "closer to horizon harder to reveal" and "spectrum collision harder to reveal" into a single complexity law: depth and separation jointly determine the minimal certificate order in logarithmic potential manner.
\end{corollary}
\begin{proof}[Proof sketch]
By Proposition \ref{prop:xi-toeplitz-negative-spectrum-margin-threshold}, we only need \(\varepsilon_N<\lambda_{\min}(\mathsf P)/2\).
Substituting the lower bound from Corollary \ref{cor:xi-pick-poisson-gap-lower-bound} for \(\lambda_{\min}(\mathsf P)\), and solving \(N\) from
\(\varepsilon_N\le C_0e^{-cNh_{\min}}\) yields a sufficient condition; constant \(C\) absorbs fixed terms like \(c^{-1}\) and \(\log 2\). Proven.
\end{proof}

\begin{proposition}[Scaling law for negative spectrum determinant stability: Poisson--Vandermonde covariance under uniform scaling]\label{prop:xi-pick-poisson-det-scaling-law}
Let \(\mathcal{C}(z)=\frac{z-\mathrm{i}}{z+\mathrm{i}}\) be the Cayley transform from upper half-plane to unit disk, and \(\mathcal{C}^{-1}(w)=\mathrm{i}\frac{1+w}{1-w}\).
For any \(m>0\), define the disk automorphism
\[
\phi_m(w):=\mathcal{C}\!\bigl(m\,\mathcal{C}^{-1}(w)\bigr),
\]
which corresponds to uniform scaling \(z\mapsto mz\) on the upper half-plane and fixes endpoint \(-1=\mathcal{C}(0)\).
Let \(a_j(m):=\phi_m(a_j)\), and let \(\mathsf P(m)\) be the Pick--Poisson Gram matrix constructed from \(\{a_j(m)\}\) via \eqref{eq:xi-pick-poisson-gram}.
Denote \(p_j(m):=P_{a_j(m)}(-1)\), \(\rho_{jk}(m):=\rho(a_j(m),a_k(m))\).
Then for all \(m>0\) there holds the covariant factorization
\[
\boxed{\;
\det\mathsf P(m)
=\Bigl(\prod_{j=1}^{\kappa}p_j(m)\Bigr)\Bigl(\prod_{1\le j<k\le\kappa}\rho_{jk}(m)^2\Bigr).\;}
\]
Moreover, \(\rho_{jk}(m)=\rho_{jk}\) (pseudohyperbolic distance is invariant under disk automorphisms), and
\[
\boxed{\;p_j(m)=m^{-1}p_j,\qquad \det\mathsf P(m)=m^{-\kappa}\det\mathsf P.\;}
\]
Therefore, the function \(s=\log m\mapsto -\log\det\mathsf P(e^s)\) is affine (hence convex), giving a deterministic dissipation direction for the scaling flow on "negative spectrum volume" (\(\det\mathsf P\)).
\end{proposition}
\begin{proof}[Proof sketch]
The first formula is the pointwise application of the same factorization in Proposition \ref{prop:xi-pick-poisson-det-factorization} on \(\{a_j(m)\}\).
Disk automorphisms preserve pseudohyperbolic distance is a classical fact, so \(\rho_{jk}(m)=\rho_{jk}\).
\par
To prove \(p_j(m)=m^{-1}p_j\), denote \(z=\mathcal{C}^{-1}(w)=x+\mathrm{i}y\).
Direct calculation gives
\[
P_{w}(-1)=\frac{1-|w|^2}{|1+w|^2}
=\frac{y}{x^2+y^2}
=\frac{\Im z}{|z|^2}.
\]
Under scaling \(z\mapsto mz\), \(\Im(mz)=m\Im z\), \(|mz|^2=m^2|z|^2\), thus
\(
P_{\phi_m(w)}(-1)=\Im(mz)/|mz|^2=m^{-1}\Im z/|z|^2
=m^{-1}P_w(-1)
\).
For \(w=a_j\) this gives \(p_j(m)=m^{-1}p_j\), thus \(\det\mathsf P(m)=m^{-\kappa}\det\mathsf P\).
The last statement follows by taking logarithm. Proven.
\end{proof}

\begin{definition}[Generalized defect entropy (negative spectrum volume potential)]\label{def:xi-generalized-defect-entropy-logdet}
Under the finite point defect setting of Proposition \ref{prop:xi-toeplitz-negative-spectrum-pick-poisson-gram}, define the Pick--Poisson negative spectrum volume potential (also called generalized defect entropy)
\[
\boxed{\;
S_{\mathrm{gen}}
:=\mathcal{E}
:=-\log\det\mathsf P\; }.
\]
\end{definition}

\begin{corollary}[Pick--Poisson negative spectrum volume "Coulomb energy" factorization: self-energy + interaction energy]\label{cor:xi-pick-poisson-energy-coulomb-decomposition}
Following the notation of \(\mathsf P\), let endpoint Poisson weights \(p_j:=P_{a_j}(-1)\), pseudohyperbolic distance \(\rho_{jk}:=\rho(a_j,a_k)\).
Then there holds the strict energy factorization
\[
\boxed{\;
\mathcal{E}
=\sum_{j=1}^{\kappa}\bigl(-\log p_j\bigr)\ +\ 2\sum_{1\le j<k\le\kappa}\bigl(-\log\rho_{jk}\bigr).\;}
\]
The first term is endpoint self-energy (single-point Busemann/Poisson cost), the second is interaction energy (pseudohyperbolic mutual interaction).
Therefore, in the semiclassical limit, the entire logarithmic potential of Toeplitz certificate negative spectrum is strictly sealed by two parts: "endpoint absorption + inter-point interaction," without need for additional continuous parameters.
\end{corollary}
\begin{proof}
By Proposition \ref{prop:xi-pick-poisson-det-factorization},
\(
\det\mathsf P=(\prod_j p_j)(\prod_{j<k}\rho_{jk}^2)
\);
taking negative logarithm of both sides yields the result. Proven.
\end{proof}

\begin{proposition}[Submultiplicativity for unions and cross-interaction terms: strict superadditivity of negative spectrum energy]\label{prop:xi-pick-poisson-union-submultiplicativity-cross-energy}
Let the defect point multiset decompose into disjoint union \(A\sqcup B\) (counted with multiplicity).
Denote corresponding Pick--Poisson Gram matrices as \(\mathsf P(A)\), \(\mathsf P(B)\), \(\mathsf P(A\cup B)\).
Then there holds the exact identity
\[
\boxed{\;
-\log\det\mathsf P(A\cup B)
=-\log\det\mathsf P(A)-\log\det\mathsf P(B)
\;+\;2\sum_{a\in A}\sum_{b\in B}\bigl(-\log\rho(a,b)\bigr).\;}
\]
Since \(\rho(a,b)\in(0,1)\), the cross term is strictly nonnegative, thus yielding the submultiplicative law
\[
\boxed{\;
\det\mathsf P(A\cup B)\ \le\ \det\mathsf P(A)\,\det\mathsf P(B).\;}
\]
Moreover, equality can only be achieved in the limiting sense (\(\rho(a,b)\to 1\) for all cross pairs).
Therefore, union energy is strictly greater than the sum of component energies, with the gap completely determined by cross pseudohyperbolic separation.
\end{proposition}
\begin{proof}
Apply the factorization in Proposition \ref{prop:xi-pick-poisson-det-factorization} to the set \(A\cup B\), and split the double product into "within-set pairs + cross pairs" to obtain the identity;
submultiplicativity follows from the cross term being nonnegative. Proven.
\end{proof}

\begin{corollary}[Extremal upper bound on negative spectrum volume: optimal achievability under fixed endpoint flux]\label{cor:xi-pick-poisson-energy-fixed-flux-extremal}
Denote total endpoint flux \(p_\Sigma:=\mathrm{tr}(\mathsf P)=\sum_{j=1}^{\kappa}p_j\).
Then there holds
\[
\boxed{\;
\det\mathsf P\ \le\ \prod_{j=1}^{\kappa}p_j\ \le\ \left(\frac{p_\Sigma}{\kappa}\right)^{\kappa},\;}
\]
thus
\[
\boxed{\;
\mathcal{E}=-\log\det\mathsf P\ \ge\ \kappa\log\kappa-\kappa\log p_\Sigma.\;}
\]
This lower bound approaches its infimum in the sense of "weights tending to uniformity and points to limiting separation," indicating that under fixed endpoint flux constraint, the optimal achievability of negative spectrum volume is jointly determined by uniformization (weights) and limiting separation (geometry).
\end{corollary}
\begin{proof}
By \(\rho_{jk}\le 1\), \(\det\mathsf P\le \prod_j p_j\).
Then by AM--GM inequality \(\prod_j p_j\le (p_\Sigma/\kappa)^\kappa\), yielding the stated upper bound;
taking negative logarithm gives the energy lower bound. Proven.
\end{proof}

\begin{corollary}[Spectral norm seal and explicit condition number bound: unified pathological fraction for near-boundary/collision]\label{cor:xi-pick-poisson-condition-number-certificate}
The matrix \(\mathsf P\) is a Gram matrix, and \(|\mathsf P_{jk}|\le \sqrt{p_jp_k}\).
Therefore, its spectral norm satisfies the uniform upper bound
\[
\boxed{\;
\lambda_{\max}(\mathsf P)\ \le\ \mathrm{tr}(\mathsf P)=p_\Sigma.\;}
\]
Combining with the spectral gap lower bound in Corollary \ref{cor:xi-pick-poisson-gap-lower-bound}, we obtain the condition number seal
\[
\boxed{\;
\kappa(\mathsf P):=\frac{\lambda_{\max}(\mathsf P)}{\lambda_{\min}(\mathsf P)}
\ \le\
\frac{p_\Sigma^{\kappa}}{\bigl(\prod_{j}p_j\bigr)\bigl(\prod_{j<k}\rho_{jk}^2\bigr)}.\;}
\]
Thus, pathology induced by "near-boundary" (\(p_j\downarrow 0\)) and "collision" (\(\rho_{jk}\downarrow 0\)) is strictly captured at the operator level by the same fractional invariant.
\end{corollary}
\begin{proof}
By positive definiteness, \(\lambda_{\max}\le \mathrm{tr}\).
Using \(\lambda_{\min}\ge \det/(\mathrm{tr})^{\kappa-1}\) (Corollary \ref{cor:xi-pick-poisson-gap-lower-bound}) and substituting the factorization of \(\det\mathsf P\) yields the result. Proven.
\end{proof}

\begin{corollary}[Logarithmic volume of Toeplitz negative spectrum: existence of limit for finite-order certificate "negative cone capacity"]\label{cor:xi-toeplitz-negative-log-volume-limit}
Assume \(\mathbf{T}_N\) enters the stable visibility interval, thus having \(\kappa\) negative eigenvalues
\(\lambda_1^{-}(\mathbf{T}_N)\le\cdots\le\lambda_{\kappa}^{-}(\mathbf{T}_N)<0\).
Define negative cone capacity (logarithmic volume)
\[
\mathcal{V}_N:=\sum_{j=1}^{\kappa}\log\bigl(-\lambda_j^{-}(\mathbf{T}_N)\bigr).
\]
Then in the semiclassical limit \(N\to\infty\) and \(N(1-|a_j|)\to\infty\) (for all \(j\)), the limit exists and satisfies
\[
\boxed{\;
\lim_{N\to\infty}\mathcal{V}_N=\log\det\mathsf P.\;}
\]
Therefore, \(\log\det\mathsf P\) is not only the limit of the product of negative eigenvalues (Proposition \ref{prop:xi-pick-poisson-det-factorization}) but also the strict limit invariant of "finite-order certificate negative subspace volume"; its degeneracy mechanism is completely classified by the energy factorization in Corollary \ref{cor:xi-pick-poisson-energy-coulomb-decomposition}.
\end{corollary}
\begin{proof}[Proof sketch]
By Proposition \ref{prop:xi-pick-poisson-det-factorization},
\(
\prod_{j=1}^{\kappa}(-\lambda_j^{-}(\mathbf{T}_N))\to \det\mathsf P
\).
On the other hand, under the margin condition of Proposition \ref{prop:xi-toeplitz-negative-spectrum-margin-threshold},
\(-\lambda_j^{-}(\mathbf{T}_N)\ge \lambda_{\min}(\mathsf P)/2>0\),
thus the logarithm is well-defined and continuous with respect to limits; taking logarithm yields the conclusion. Proven.
\end{proof}

\begin{proposition}[Spectral stability under boundary density perturbation: negative spectrum eigenvalues Lipschitz under \texorpdfstring{$L^1$}{L1} perturbation in separation domain]\label{prop:xi-boundary-density-perturbation-spectral-stability}
Let basepoint scanning profiles \(H\) be as defined in Definition \ref{def:xi-basepoint-scan-profile}, and assume two groups of finite point defects induce \(H,\widetilde H\).
Assume both parameter groups are in the separation domain: there exist \(\rho_0>0\), \(p_0>0\) such that
\(
\rho_{\min}\ge \rho_0
\)
and
\(
\sum_j p_j\ge p_0
\)
(for both groups).
Then there exists a constant \(C=C(\kappa,\rho_0,p_0)\) such that the corresponding Pick--Poisson matrices satisfy
\[
\boxed{\;
\|\mathsf P-\widetilde{\mathsf P}\|_{\mathrm{op}}
\ \le\ C\,\|H-\widetilde H\|_{L^1(\RR)}.\;}
\]
Thus in the stable visibility interval, Toeplitz negative eigenvalues satisfy a Weyl-type stability bound
\[
\boxed{\;
\max_{1\le j\le\kappa}\bigl|\lambda_j^{-}(\mathbf{T}_N)-\widetilde\lambda_j^{-}(\widetilde{\mathbf{T}}_N)\bigr|
\ \le\ C\,\|H-\widetilde H\|_{L^1(\RR)}\ +\ o_N(1),\;}
\]
where \(o_N(1)\to 0\) is the certificate truncation error (controlled by \(\varepsilon_N\)).
This conclusion strictly linearizes "boundary visible perturbation" and "certificate negative spectrum perturbation" in the separation domain, giving an auditable robustness seal.
\end{proposition}
\begin{proof}[Proof sketch]
In the separation domain with fixed \(\kappa\) and away from near-boundary/collision degeneracy, the map \(\{a_j\}\mapsto H\mapsto \mathsf P\) is smooth and has uniformly bounded derivatives on compact domains;
by compactness and mean value theorem, we obtain Lipschitz constant control for \(L^1\) perturbation, thus deriving the linear upper bound for \(\|\mathsf P-\widetilde{\mathsf P}\|_{\mathrm{op}}\).
For Toeplitz negative spectrum, first use Proposition \ref{prop:xi-toeplitz-negative-spectrum-pick-poisson-gram} to align its restriction on the negative subspace with \(-\mathsf P\), then use Weyl's inequality and truncation error \(\varepsilon_N=o(1)\) to obtain the second formula. Proven.
\end{proof}

\paragraph{Image charge mirror energy and weighted Fekete variation: upper half-plane realization of Pick-Poisson negative spectral volume potential}
This paragraph elevates generalized defect entropy of Definition \ref{def:xi-generalized-defect-entropy-logdet}
\(
S_{\mathrm{gen}}=-\log\det\mathsf P
\)
to discretization of upper half-plane Dirichlet Green energy, thereby providing a suite of independently citable electrostatic characterizations including balance equations, stabilizer subgroup covariance laws, and variational principles.

\begin{proposition}[Image charge mirror energy representation: \texorpdfstring{$-\log\det\mathsf P$}{-log det P} equals discretization of upper half-plane Dirichlet Green energy]\label{prop:xi-pick-poisson-image-charge-green-energy}
Let finite point-wise defect points be \(\{a_j\}_{j=1}^{\kappa}\subset\mathbb{D}\) (expanded by multiplicity), and let \(\mathsf P\) be their Pick-Poisson Gram matrix \eqref{eq:xi-pick-poisson-gram}.
Define "endpoint-normalized" upper half-plane coordinates
\[
z_j:=\mathcal{C}^{-1}(a_j)\in\mathbb{H},
\qquad
t_j:=-\frac{1}{z_j}\in\mathbb{H},
\qquad
t_j=\gamma_j+\mathrm{i}y_j.
\]
Then endpoint Poisson weight satisfies
\[
\boxed{\;p_j:=P_{a_j}(-1)=y_j=\Im(t_j).\;}
\]
And pseudo-hyperbolic distance satisfies
\[
\boxed{\;
\rho(a_j,a_k)=\left|\frac{t_j-t_k}{t_j-\overline t_k}\right|\qquad(1\le j,k\le\kappa).\;}
\]
Thus generalized defect entropy (negative spectral volume potential)
\(
S_{\mathrm{gen}}:=-\log\det\mathsf P
\)
has strict equivalent form
\[
\boxed{\;
S_{\mathrm{gen}}
=\sum_{j=1}^{\kappa}\bigl(-\log y_j\bigr)
\;+\;
2\sum_{1\le j<k\le\kappa}\Bigl(\log|t_j-\overline t_k|-\log|t_j-t_k|\Bigr).\;}
\]
Let upper half-plane Dirichlet Green function be
\[
G_{\mathbb{H}}(t,s):=\log\left|\frac{t-\overline s}{t-s}\right|,
\]
then above formula can be written as
\[
S_{\mathrm{gen}}
=\sum_{j=1}^{\kappa}\bigl(-\log\Im(t_j)\bigr)
\;+\;
2\sum_{1\le j<k\le\kappa}G_{\mathbb{H}}(t_j,t_k).
\]
Therefore "compression energy" of point-wise defect under endpoint horizon has strict electrostatic realization: external energy completely closed by Green mutual interaction of real points with their boundary mirror images, while self-energy term \(-\sum_j\log\Im(t_j)\) provides non-eliminable cost of endpoint absorption.
\end{proposition}
\begin{proof}[Proof outline]
From computation in Proposition \ref{prop:xi-pick-poisson-det-scaling-law},
\(
P_w(-1)=\Im z/|z|^2
\) (where \(z=\mathcal{C}^{-1}(w)\)).
Taking \(w=a_j\) and using \(\Im(-1/z)=\Im z/|z|^2\) yields \(p_j=\Im(t_j)=y_j\).
\par
Pseudo-hyperbolic distance under Cayley transformation satisfies \(\rho(a_j,a_k)=\left|\frac{z_j-z_k}{z_j-\overline z_k}\right|\), while mapping \(z\mapsto -1/z\) belongs to \(\mathrm{Aut}(\mathbb{H})\) preserving this invariant, hence also \(\rho(a_j,a_k)=\left|\frac{t_j-t_k}{t_j-\overline t_k}\right|\).
Substituting \(p_j\) and \(\rho_{jk}\) into factorization of Proposition \ref{prop:xi-pick-poisson-det-factorization} and taking negative logarithm yields energy representation.
\end{proof}

\begin{proposition}[Balance equations: critical points of \texorpdfstring{$S_{\mathrm{gen}}$}{S_gen} satisfy strict cancellation of mirror Coulomb forces]\label{prop:xi-pick-poisson-energy-balance-equations}
Under upper half-plane parameters \((\gamma_j,y_j)\) of Proposition \ref{prop:xi-pick-poisson-image-charge-green-energy}, view \(S_{\mathrm{gen}}\) as function of variables.
Then its Euclidean gradient has explicit "mirror force difference" form: for any \(j\),
\[
\boxed{\;
\frac{\partial S_{\mathrm{gen}}}{\partial \gamma_j}
=
2\sum_{k\neq j}\left(
\frac{\gamma_j-\gamma_k}{(\gamma_j-\gamma_k)^2+(y_j+y_k)^2}
-
\frac{\gamma_j-\gamma_k}{(\gamma_j-\gamma_k)^2+(y_j-y_k)^2}
\right)\;}
\]
and
\[
\boxed{\;
\frac{\partial S_{\mathrm{gen}}}{\partial y_j}
=
-\frac{1}{y_j}
+
2\sum_{k\neq j}\left(
\frac{y_j+y_k}{(\gamma_j-\gamma_k)^2+(y_j+y_k)^2}
-
\frac{y_j-y_k}{(\gamma_j-\gamma_k)^2+(y_j-y_k)^2}
\right).\;}
\]
Therefore any interior extremal configuration must satisfy above \(\kappa\) equations simultaneously being zero; their meaning is: each point's tangential and normal Coulomb forces under boundary conductor (mirror) action both reach strict balance, while self-energy term \(-\log y_j\) only provides non-eliminable endpoint attraction in normal equation.
\end{proposition}
\begin{proof}[Proof outline]
From energy expansion in Proposition \ref{prop:xi-pick-poisson-image-charge-green-energy},
\[
S_{\mathrm{gen}}
=-\sum_j\log y_j
\;+\;\sum_{j<k}\log\!\bigl((\gamma_j-\gamma_k)^2+(y_j+y_k)^2\bigr)
\;-\;\sum_{j<k}\log\!\bigl((\gamma_j-\gamma_k)^2+(y_j-y_k)^2\bigr).
\]
Taking partial derivatives with respect to \(\gamma_j\), \(y_j\) term-wise and arranging symmetric terms yields stated expressions.
\end{proof}

\begin{corollary}[Horizon stabilizer subgroup covariance law: translation invariant, scaling has \texorpdfstring{$\kappa\log\lambda$}{k log lambda} cofactor]\label{cor:xi-pick-poisson-energy-stabilizer-covariance}
Let upper half-plane point set undergo horizontal translation \(t_j\mapsto t_j+x_0\) (\(x_0\in\RR\)) and positive scaling \(t_j\mapsto \lambda t_j\) (\(\lambda>0\)).
Then
\[
\boxed{\;
S_{\mathrm{gen}}(t_1+x_0,\dots,t_\kappa+x_0)=S_{\mathrm{gen}}(t_1,\dots,t_\kappa),\;}
\]
and
\[
\boxed{\;
S_{\mathrm{gen}}(\lambda t_1,\dots,\lambda t_\kappa)=S_{\mathrm{gen}}(t_1,\dots,t_\kappa)-\kappa\log\lambda.\;}
\]
In other words, inter-point mutual energy term (Green difference) is strictly invariant under horizon stabilizer subgroup, while self-energy term produces unique computable homology cofactor \(-\kappa\log\lambda\) through \(-\sum_j\log\Im(t_j)\).
\end{corollary}
\begin{proof}
Translation preserves differences \(t_j-t_k\), \(t_j-\overline t_k\) and preserves \(\Im(t_j)\), hence energy invariant.
Under scaling \(\Im(\lambda t_j)=\lambda\Im(t_j)\), while \(\log|\lambda(t_j-\overline t_k)|-\log|\lambda(t_j-t_k)|\)'s \(\log\lambda\) terms cancel, hence only self-energy term contributes \(-\kappa\log\lambda\).
\end{proof}

\begin{proposition}[Weighted Fekete variational principle: maximizing \texorpdfstring{$\det\mathsf P$}{det P} equivalent to minimizing mirror energy]\label{prop:xi-pick-poisson-weighted-fekete-variational}
Fix \(\kappa\) and take constraint set \(\mathcal{K}\subset\mathbb{H}^{\kappa}\).
View \(S_{\mathrm{gen}}=-\log\det\mathsf P\) as extended real-valued function: if \(\det\mathsf P=0\) take \(S_{\mathrm{gen}}=+\infty\).
Then extremal problems
\[
\max_{(t_1,\dots,t_\kappa)\in\mathcal{K}}\ \det\mathsf P(t_1,\dots,t_\kappa)
\qquad\Longleftrightarrow\qquad
\min_{(t_1,\dots,t_\kappa)\in\mathcal{K}}\ S_{\mathrm{gen}}(t_1,\dots,t_\kappa)
\]
are strictly equivalent.
When \(\mathcal{K}\) is compact and contains at least one configuration making \(\det\mathsf P>0\), above maximum and minimum both exist and are finite.
If further imposing geometric separation lower bound \(\rho_{\min}\ge\rho_0>0\) and height lower bound \(y_{\min}>0\), then extremal sequences are precompact, and any interior extremal point must satisfy balance equations of Proposition \ref{prop:xi-pick-poisson-energy-balance-equations}.
\end{proposition}
\begin{proof}[Proof outline]
Equivalence comes from monotonicity of logarithm:
\(
\arg\max\det\mathsf P=\arg\min(-\log\det\mathsf P)
\).
If \(\mathcal{K}\) compact, then \(\det\mathsf P\) is continuous function, hence maximum exists.
On the other hand, if \(\mathcal{K}\) contains point making \(\det\mathsf P>0\), then \(\inf_{\mathcal{K}}S_{\mathrm{gen}}<\infty\);
and \(S_{\mathrm{gen}}\) as extended real-valued function of \(-\log\circ\det\) is lower semicontinuous on \(\mathcal{K}\), hence minimum exists and is finite.
\par
Separation/height lower bounds exclude edge-touching and collision degeneracies (thus excluding energy divergence and ensuring \(\det\mathsf P\) uniformly bounded away from zero in this domain), giving precompactness; interior extremal point satisfies gradient zero yields balance equations.
\end{proof}

\begin{proposition}[Determinant chain decomposition: stepwise ledger law of conditional flux]\label{prop:xi-pick-poisson-schur-flux-ledger}
For any fixed point order \(1,\dots,\kappa\), denote \(\mathsf P^{(m)}\) as principal submatrix induced by first \(m\) points.
Define Schur complement (conditional flux)
\[
\pi_m
:=
\mathsf P_{mm}-\mathsf P_{m,1:m-1}\,(\mathsf P^{(m-1)})^{-1}\,\mathsf P_{1:m-1,m},
\qquad
m=1,\dots,\kappa,
\]
where \(\mathsf P^{(0)}\) is conventionally empty.
Then there holds strict product law
\[
\boxed{\;
\det\mathsf P=\prod_{m=1}^{\kappa}\pi_m,
\qquad
S_{\mathrm{gen}}=-\sum_{m=1}^{\kappa}\log\pi_m.\;}
\]
And each \(\pi_m\) satisfies
\[
\boxed{\;
0<\pi_m\le \mathsf P_{mm}=p_m.\;}
\]
Therefore \(S_{\mathrm{gen}}\) has strict "stepwise externalization ledger" structure: each time a defect is added, new energy exactly equals its conditional flux loss on space spanned by existing defects.
\end{proposition}
\begin{proof}[Proof outline]
\(\det\mathsf P=\prod_m\pi_m\) is standard conclusion of Schur complement chain decomposition (equivalent to Cholesky decomposition) of positive definite matrix.
Upper bound \(\pi_m\le \mathsf P_{mm}\) can be interpreted via Gram-Schmidt: when \(\mathsf P\) is Gram matrix of vector family, \(\pi_m\) equals squared norm of \(m\)-th vector on orthogonal complement of previously spanned subspace, hence not exceeding its full norm squared \(\mathsf P_{mm}\).
\end{proof}

\begin{proposition}[Boundary interpolation margin interpretation: \texorpdfstring{$\det\mathsf P$}{det P} and \texorpdfstring{$\lambda_{\min}(\mathsf P)$}{lambda_min(P)} as verifiable margin certificates]\label{prop:xi-pick-poisson-interpolation-margin-certificate}
Let \(k_a(w)=(1-\overline a\,w)^{-1}\) be Hardy reproducing kernel, and let
\[
\psi_j(w):=\frac{1-|a_j|^2}{1+\overline a_j}\,k_{a_j}(w)
=\frac{1-|a_j|^2}{(1+\overline a_j)(1-\overline a_j w)}\in H^2.
\]
Then \(\mathsf P\) can be written as Gram matrix of \(\{\psi_j\}\):
\[
\boxed{\;
\langle \psi_j,\psi_k\rangle_{H^2}=\mathsf P_{jk}.\;}
\]
Therefore \(\det\mathsf P\) equals squared volume of parallelepiped spanned by \(\{\psi_j\}\) in \(H^2\) (Gram determinant theorem), with non-degeneracy equivalent to linear independence of endpoint-anchored evaluation subspace.
And
\[
\lambda_{\min}(\mathsf P)
=\min_{\|c\|_2=1}\left\|\sum_{j=1}^{\kappa}c_j\psi_j\right\|_{H^2}^2
\]
provides worst "margin" of this evaluation subspace.
Combining explicit lower bound of Corollary \ref{cor:xi-pick-poisson-gap-lower-bound} yields fully geometrized verifiable margin certificate: its degeneration mechanism strictly attributed to edge-touching (\(p_j\downarrow 0\)) and collision (\(\rho_{jk}\downarrow 0\)).
\end{proposition}
\begin{proof}[Proof outline]
From \(\langle k_{a_j},k_{a_k}\rangle_{H^2}=(1-\overline a_j a_k)^{-1}\) and
\(
\psi_j=\frac{1-|a_j|^2}{1+\overline a_j}k_{a_j}
\),
we get
\[
\langle \psi_j,\psi_k\rangle_{H^2}
=\frac{(1-|a_j|^2)(1-|a_k|^2)}{(1+\overline a_j)(1+a_k)(1-\overline a_j a_k)}
=\mathsf P_{jk}.
\]
Gram determinant and minimal eigenvalue variational representation are standard linear algebra facts.
\end{proof}

\begin{definition}[Anchored capacity (discrete horizon capacity)]\label{def:xi-anchored-capacity-capminus1}
Let \(A=\{a_j\}_{j=1}^{\kappa}\subset\mathbb{D}\) be finite point-wise defect point set (expanded by multiplicity), and let \(\mathsf P(A)\) be its Pick-Poisson Gram matrix \eqref{eq:xi-pick-poisson-gram}.
Define anchored capacity (discrete horizon capacity)
\[
\boxed{\;
\mathrm{cap}_{-1}(A):=\bigl(\det\mathsf P(A)\bigr)^{\frac{1}{2\kappa}}
=\exp\!\left(-\frac{1}{2\kappa}S_{\mathrm{gen}}(A)\right)\ \in(0,\infty).\;}
\]
\end{definition}

\begin{proposition}[Three properties of anchored capacity: permutation invariance, degeneration criterion, and union submultiplicativity]\label{prop:xi-anchored-capacity-three-properties}
Anchored capacity \(\mathrm{cap}_{-1}\) has following properties:
\begin{enumerate}
  \item \textbf{(Permutation invariance)} For any rearrangement \(A\mapsto A'\) (preserving multiset) we have \(\mathrm{cap}_{-1}(A')=\mathrm{cap}_{-1}(A)\);
  \item \textbf{(Monotone degeneration)} If there exists sequence of configurations such that some endpoint weight \(p_j=P_{a_j}(-1)\downarrow 0\) (remaining points stay in compact domain), or collision exists \(\rho(a_j,a_k)\downarrow 0\) (remaining points fixed), then \(\mathrm{cap}_{-1}(A)\to 0\);
  \item \textbf{(Union submultiplicativity)} For disjoint union \(A\sqcup B\) (counting multiplicity)
  \[
  \det\mathsf P(A\cup B)\ \le\ \det\mathsf P(A)\,\det\mathsf P(B),
  \]
  thus
  \[
  \boxed{\;
  \mathrm{cap}_{-1}(A\cup B)
  \le
  \mathrm{cap}_{-1}(A)^{\frac{|A|}{|A|+|B|}}\,
  \mathrm{cap}_{-1}(B)^{\frac{|B|}{|A|+|B|}}.\;}
  \]
\end{enumerate}
Therefore \(\mathrm{cap}_{-1}\) at regime level provides single-scalar scale of "point-wise defect can be stably resolved by external interface"; its logarithmic potential \(S_{\mathrm{gen}}=-\log\det\mathsf P\) is additive ledger potential of discrete capacity.
\end{proposition}
\begin{proof}[Proof outline]
(1) Follows from determinant invariance under index permutation.
\par
(2) Follows from determinant factorization \(\det\mathsf P=(\prod_j p_j)(\prod_{j<k}\rho_{jk}^2)\) (Proposition \ref{prop:xi-pick-poisson-det-factorization}).
\par
(3) From union cross-energy identity (Proposition \ref{prop:xi-pick-poisson-union-submultiplicativity-cross-energy})
\(
\det\mathsf P(A\cup B)=\det\mathsf P(A)\det\mathsf P(B)\cdot \prod_{a\in A,b\in B}\rho(a,b)^2
\le \det\mathsf P(A)\det\mathsf P(B)
\);
then take \(1/(2(|A|+|B|))\) power.
\end{proof}

\begin{proposition}[Anchored Hardy minimal norm interpolation: \texorpdfstring{$\mathsf P\succ 0$}{P>0} if and only if anchored interpolation feasible]\label{prop:xi-anchored-hardy-interpolation-pick-equivalence}
Continue anchored kernel vectors \(\psi_j\in H^2\) of Proposition \ref{prop:xi-pick-poisson-interpolation-margin-certificate}, and define anchored evaluation operator
\[
\mathcal{E}_{A,-1}:H^2\to\CC^{\kappa},
\qquad
(\mathcal{E}_{A,-1}f)_j:=\langle f,\psi_j\rangle_{H^2}
=\frac{1-|a_j|^2}{1+\overline a_j}\,f(a_j).
\]
Given data \(c\in\CC^{\kappa}\), consider anchored interpolation problem: minimize \(\|f\|_{H^2}\) under constraint \(\mathcal{E}_{A,-1}f=c\).
Then following conclusions hold:
\begin{enumerate}
  \item \(\mathsf P(A)\succ 0\) if and only if for each \(c\) there exists unique minimal norm solution \(f_c\in H^2\);
  \item When \(\mathsf P(A)\succ 0\), minimal norm satisfies closed form
  \[
  \boxed{\;
  \|f_c\|_{H^2}^2=c^\ast\,\mathsf P(A)^{-1}\,c;\;}
  \]
  \item Hence stability constant of interpolation operator \(c\mapsto f_c\) is completely sealed by \(\lambda_{\min}(\mathsf P(A))\):
  \[
  \boxed{\;
  \|f_c\|_{H^2}^2\le \lambda_{\min}(\mathsf P(A))^{-1}\,\|c\|_2^2,}
  \]
  with this upper bound attainable in extremal direction (taking \(c\) as minimal eigenvector).
\end{enumerate}
Therefore \(\mathsf P\) is simultaneously Gram matrix of "defect geometry" and Pick matrix of "anchored minimal norm interpolation"; \(\det\mathsf P\) and \(\lambda_{\min}(\mathsf P)\) respectively characterize volume-type and margin-type verifiable margins.
\end{proposition}
\begin{proof}[Proof outline]
\(\mathsf P(A)\succ 0\) equivalent to \(\{\psi_j\}\) linearly independent, thus operator \(\mathcal{E}_{A,-1}\) has full rank on \(\Span\{\psi_j\}\).
Minimal norm problem in Hilbert space satisfies representer theorem: solution necessarily has form \(f_c=\sum_j \alpha_j\psi_j\).
Substituting into constraint yields linear equation \(\mathsf P(A)\alpha=c\), hence unique solution is \(\alpha=\mathsf P(A)^{-1}c\).
Thus
\(
\|f_c\|^2=\alpha^\ast \mathsf P(A)\alpha=c^\ast\mathsf P(A)^{-1}c
\).
Stability bound follows from Rayleigh quotient estimate immediately.
\end{proof}

\begin{proposition}[Minimal certificate order interpolation interpretation: Toeplitz development threshold equivalent to anchored interpolation condition number threshold]\label{prop:xi-toeplitz-threshold-via-anchored-interpolation}
Let \(A=\{a_j\}_{j=1}^{\kappa}\subset\mathbb{D}\) in separation domain and let \(h_{\min}:=\min_j(1-|a_j|)\).
Let \(f_c\) be minimal norm solution of Proposition \ref{prop:xi-anchored-hardy-interpolation-pick-equivalence}, and let \(\Pi_N:H^2\to\Span\{1,w,\dots,w^N\}\) be degree truncation orthogonal projection, denote \(f_c^{(N)}:=\Pi_N f_c\).
Then there exists constant \(C>0\) (depending at most on \(\kappa\)) such that: if
\[
\boxed{\;
N\ \ge\ C\,\frac{\log\kappa(\mathsf P(A))}{h_{\min}},\qquad
\kappa(\mathsf P):=\frac{\lambda_{\max}(\mathsf P)}{\lambda_{\min}(\mathsf P)},\;}
\]
then for all \(c\in\CC^{\kappa}\)
\[
\boxed{\;
\|f_c-f_c^{(N)}\|_{H^2}\ \le\ \frac12\,\|f_c\|_{H^2}.\;}
\]
Thus beyond this threshold, negative spectral margin of finite-order Toeplitz certificate on defect kernel spanned subspace is same order as spectral margin of \(\mathsf P\) and stably develops (consistent with threshold laws of Proposition \ref{prop:xi-toeplitz-negative-spectrum-margin-threshold} and Corollary \ref{cor:xi-toeplitz-visibility-complexity-threshold}).
Therefore "insufficient certificate order causing defect invisible" in strict sense is equivalent to "anchored interpolation condition number too large causing truncation uncontrollable".
\end{proposition}
\begin{proof}[Proof outline]
From \(k_{a}(w)=\sum_{n\ge 0}\overline a^{\,n}w^n\) get \(\|\Pi_{>N}k_a\|_{H^2}=|a|^{N+1}(1-|a|^2)^{-1/2}\).
For anchored vector \(\psi_j=\frac{1-|a_j|^2}{1+\overline a_j}k_{a_j}\) get
\(
\|\Pi_{>N}\psi_j\|_{H^2}\le \sqrt{p_j}\,|a_j|^{N+1}\le \sqrt{p_j}\,e^{-Nh_{\min}}
\).
Also \(f_c=\sum_j\alpha_j\psi_j\) with \(\alpha=\mathsf P^{-1}c\), thus \(\|\alpha\|_2\le \lambda_{\min}(\mathsf P)^{-1}\|c\|_2\).
Merging tail estimate with \(\kappa(\mathsf P)=\lambda_{\max}/\lambda_{\min}\) yields stated "for all \(c\)" relative error threshold; constant \(C\) absorbs \(\kappa\) dimension factor and upper bound of \(p_\Sigma=\mathrm{tr}\,\mathsf P\) (Corollary \ref{cor:xi-pick-poisson-condition-number-certificate}).
\end{proof}

\begin{proposition}[Deterministic \texorpdfstring{$\kappa$}{k}-DPP of point-wise defect: minimal repulsion distribution with density \texorpdfstring{$\det\mathsf P$}{det P}]\label{prop:xi-pick-poisson-kdpp}
Take finite candidate grid \(\Omega=\{x_1,\dots,x_M\}\subset\mathbb{D}\) (or its upper half-plane image).
Define Pick-Poisson positive definite kernel
\[
K_{-1}(w,z):=\frac{(1-|w|^2)(1-|z|^2)}{(1+\overline w)(1+z)(1-\overline w z)}\qquad(w,z\in\mathbb{D}),
\]
and let \(L\in\Mat_M(\CC)_{\mathrm{Herm}}\) be its Gram on \(\Omega\):
\[
L_{uv}:=K_{-1}(x_u,x_v)\qquad(1\le u,v\le M).
\]
Then \(L\succeq 0\), and for any \(S\subseteq\Omega\) we have \(\det(L_S)=\det\mathsf P(S)\).
For fixed point count \(\kappa\) and parameter \(\beta>0\), define distribution on \(\kappa\)-subsets of \(\Omega\)
\[
\mathbb{P}_{\beta}^{(\Omega)}(S)\ \propto\ \det(L_S)^{\beta}\qquad(|S|=\kappa).
\]
Then when \(\beta=1\), above distribution is standard \(\kappa\)-DPP (fixed-size determinantal distribution; also \(k\)-DPP/\(L\)-ensemble), and can be viewed as fixed-size version obtained by conditioning \(L\)-ensemble DPP on point count being \(\kappa\) (see \cite{KuleszaTaskar2012DPP}).
And its log-density equals \(-S_{\mathrm{gen}}(S)\) (up to normalization constant depending only on \(\Omega\)), thus giving deterministic repulsion realization of "maximizing verifiable margin \(\Longleftrightarrow\) minimizing mirror energy".
\end{proposition}
\begin{proof}[Proof outline]
For finite \(\Omega\), matrix \(L\succeq 0\) defines a \(k\)-DPP, whose weight on size-\(\kappa\) subset is proportional to \(\det(L_S)\) (\cite{KuleszaTaskar2012DPP}).
And \(S_{\mathrm{gen}}(S)=-\log\det\mathsf P(S)=-\log\det(L_S)\) (Definition \ref{def:xi-generalized-defect-entropy-logdet}), hence log-density representation.
\end{proof}

\begin{proposition}[Mean field limit template: main order of \texorpdfstring{$\kappa^{-2}S_{\mathrm{gen}}$}{k^{-2}S_gen} determined by mirror Green energy functional]\label{prop:xi-pick-poisson-meanfield-template}
Under upper half-plane coordinates of Proposition \ref{prop:xi-pick-poisson-image-charge-green-energy}, let \(t_1,\dots,t_\kappa\in\mathcal{K}\subset\mathbb{H}\) (\(\mathcal{K}\) compact with \(\Im t\ge y_0>0\)).
Suppose empirical measure \(\mu_\kappa:=\frac{1}{\kappa}\sum_{j=1}^{\kappa}\delta_{t_j}\) weakly converges to \(\mu\in\mathcal{P}(\mathcal{K})\), and point configuration does not undergo macroscopic collapse at scale (e.g., uniform separation lower bound \(\rho_{\min}\ge\rho_0>0\) exists).
Then there holds main order limit
\[
\boxed{\;
\lim_{\kappa\to\infty}\frac{1}{\kappa^2}S_{\mathrm{gen}}(t_1,\dots,t_\kappa)
=\iint_{\mathcal{K}\times\mathcal{K}}G_{\mathbb{H}}(t,s)\,\dd\mu(t)\dd\mu(s),\;}
\]
where \(G_{\mathbb{H}}(t,s)=\log\left|\frac{t-\overline s}{t-s}\right|\).
Self-energy term \(\kappa^{-2}\sum_j(-\log\Im t_j)\) within fixed \(\mathcal{K}\) is sub-order and degenerates to \(0\).
\end{proposition}
\begin{proof}[Proof outline]
In separation domain \(G_{\mathbb{H}}\) is bounded continuous on \(\mathcal{K}\times\mathcal{K}\), hence
\(
\frac{2}{\kappa^2}\sum_{j<k}G_{\mathbb{H}}(t_j,t_k)
\to \iint G_{\mathbb{H}}\,\dd\mu\dd\mu
\);
while self-energy term is \(O(\kappa)\) hence dividing by \(\kappa^2\) vanishes.
\end{proof}

\begin{corollary}[Euler-Lagrange equation of equilibrium measure: boundary constant potential condition induced by horizon anchoring]\label{cor:xi-pick-poisson-equilibrium-euler-lagrange}
On continuous energy functional of Proposition \ref{prop:xi-pick-poisson-meanfield-template}
\(
\mathcal{E}(\mu):=\iint G_{\mathbb{H}}(t,s)\,\dd\mu(t)\dd\mu(s)
\),
if \(\mu^\ast\) is minimal equilibrium measure on \(\mathcal{P}(\mathcal{K})\), then there exists constant \(C\) such that its potential function
\[
U^{\mu^\ast}(t):=\int_{\mathcal{K}}G_{\mathbb{H}}(t,s)\,\dd\mu^\ast(s)
\]
satisfies Euler-Lagrange condition
\[
\boxed{\;
U^{\mu^\ast}(t)=C\ \text{on}\ \mathrm{supp}(\mu^\ast),\qquad
U^{\mu^\ast}(t)\ge C\ \text{on}\ \mathcal{K}.\;}
\]
This constant potential condition provides variational characterization of "most visible defect distribution" in continuous limit (see Frostman principle in logarithmic potential theory, \cite{SaffTotik1997LogarithmicPotentials}).
\end{corollary}
\begin{proof}[Proof outline]
This is standard equilibrium measure conclusion of logarithmic potential energy on compact set: minimality equivalent to potential function satisfying constant potential/barrier condition.
\end{proof}

\begin{proposition}[Concentration inequality of negative spectral volume: Lipschitz property and McDiarmid-type concentration of \texorpdfstring{$\log\det\mathsf P$}{log det P} in separation domain]\label{prop:xi-pick-poisson-logdet-concentration}
Suppose point configuration satisfies uniform separation lower bound \(\rho_{\min}\ge\rho_0>0\) and height lower bound \(\Im t_j\ge y_0>0\), and let
\(
F(t_1,\dots,t_\kappa):=\log\det\mathsf P(t_1,\dots,t_\kappa)
\).
Then \(F\) is globally Lipschitz in this separation domain (constant depending only on \(\kappa,\rho_0,y_0\)).
Therefore for independent sampling \(t_1,\dots,t_\kappa\) (or Markov chain sampling satisfying bounded difference condition) has McDiarmid-type concentration: there exists constant \(c>0\) such that for any \(u>0\),
\[
\boxed{\;
\mathbb{P}\bigl(|F-\mathbb{E}F|\ge u\bigr)\ \le\ 2\exp\!\left(-c\,\frac{u^2}{\kappa}\right)\;}
\]
(see \cite{McDiarmid1989BoundedDifferences}).
\end{proposition}
\begin{proof}[Proof outline]
From \(F=-S_{\mathrm{gen}}\) and Proposition \ref{prop:xi-pick-poisson-energy-balance-equations} providing explicit rational form of \(\nabla S_{\mathrm{gen}}\); in separation domain denominators uniformly bounded away from zero, hence gradient uniformly bounded, thus Lipschitz.
Concentration inequality under bounded difference/independent sampling follows from McDiarmid method.
\end{proof}

\begin{proposition}[Free energy dissipation under horizon coarse-graining: strict contraction of relative entropy by Poisson semigroup]\label{prop:xi-pick-poisson-poisson-semigroup-free-energy-dissipation}
Let \(T_t=\exp(-t|D|)\) be Poisson semigroup on \(\RR\) (Proposition \ref{con:discussion-dtn-generator}), and let its tensorized semigroup on \(\RR^\kappa\) be \(T_t^{\otimes \kappa}\).
For any two probability measures \(\mu,\nu\) (defined on \(\RR^\kappa\)), denote
\[
\mu_t:=\mu\,T_t^{\otimes\kappa},\qquad \nu_t:=\nu\,T_t^{\otimes\kappa},
\]
i.e., Markov pushforward of convolution with Poisson kernel on each coordinate respectively.
Then relative entropy satisfies semigroup monotonicity
\[
\boxed{\;
D(\mu_t\|\nu_t)\ \le\ D(\mu_s\|\nu_s)\qquad(t\ge s\ge 0).\;}
\]
If further choosing some reference measure \(\nu\) invariant under \(T_t^{\otimes\kappa}\) (e.g., take Haar/uniform benchmark under periodization/finite volume caliber), then \(t\mapsto D(\mu_t\|\nu)\) is monotone non-increasing;
and under smoothness conditions its derivative is given by Dirichlet form of generator \(\sum_{j=1}^{\kappa}|D_{\gamma_j}|\):
\[
\frac{\dd}{\dd t}D(\mu_t\|\nu)=-\mathcal{I}(\mu_t)\le 0,
\]
where \(\mathcal{I}\) is corresponding non-local Fisher information.
Therefore "horizon recycling/coarse-graining" still manifests as strict thermodynamic potential dissipation at multi-defect statistical ensemble level: coarse-graining is not approximation operation, but irreversible flow uniquely determined by geometric generator.
\end{proposition}
\begin{proof}[Proof outline]
From semigroup property \(\mu_t=(\mu_s)T_{t-s}^{\otimes\kappa}\), \(\nu_t=(\nu_s)T_{t-s}^{\otimes\kappa}\),
and relative entropy contraction under Markov channel (data processing inequality) get
\(
D(\mu_t\|\nu_t)\le D(\mu_s\|\nu_s)
\).
When \(\nu\) is invariant measure, \(\nu_t=\nu\) yields monotonicity of \(D(\mu_t\|\nu)\); derivative formula is standard entropy dissipation identity (expressed by generator).
\end{proof}

\begin{conjecture}[Energy-index consistency of two-source interface: order isomorphism of delinking gap and cross mutual energy]\label{conj:xi-energy-index-gap-comparability}
Let Jones commutator scalar of two-source interface be \(\lambda=\exp(-\Gamma)\) (Proposition \ref{prop:op-algebra-jones-scalar-equals-exp-minus-gap}), and denote index gap as \(\Gamma\ge 0\).
For geometric union \(A\sqcup B\) define cross mutual energy
\[
\mathcal{E}_{\times}(A,B):=2\sum_{a\in A}\sum_{b\in B}\bigl(-\log\rho(a,b)\bigr)
=-\log\frac{\det\mathsf P(A\cup B)}{\det\mathsf P(A)\det\mathsf P(B)}.
\]
In stable separation domain (\(\rho(a,b)\) bounded away from \(0\) and endpoint flux non-collapsing), there exist constants \(c_1,c_2>0\) such that
\[
\boxed{\;
c_1\,\Gamma(A,B)\ \le\ \mathcal{E}_{\times}(A,B)\ \le\ c_2\,\Gamma(A,B),\;}
\]
where \(\Gamma(A,B)\) represents index gap defined by corresponding Jones projection commutator scalar.
This alignment conjecture suggests: in stable domain, "delinking required resource" (index gap) and "union mutual energy" (negative spectral cross term) are order-isomorphic, both controlled by same separation parameter.
\end{conjecture}

\begin{proposition}[Operator Norm Formulation of Anchored Capacity: \texorpdfstring{$\mathrm{cap}_{-1}$}{cap-1} as the Geometric Mean of the Interpolation Right-Inverse Exterior Power Norm]\label{prop:xi-anchored-capacity-exterior-power-interpolant}
Following the anchored evaluation operator
\(\mathcal{E}_{A,-1}:H^2\to\CC^{\kappa}\) and the Gram matrix \(\mathsf P(A)\succ0\) from Proposition \ref{prop:xi-anchored-hardy-interpolation-pick-equivalence}.
Denote by the minimum norm principle the \emph{anchored interpolation right inverse} (minimum norm interpolation operator)
\[
\mathcal{I}_{A,-1}:\CC^{\kappa}\to H^2,\qquad \mathcal{I}_{A,-1}(c):=f_c,
\]
where \(f_c\) is the unique minimum norm solution satisfying \(\mathcal{E}_{A,-1}f=c\).
Then \(\mathcal{I}_{A,-1}\) is linear and bounded, and satisfies
\[
\boxed{\;
\mathcal{I}_{A,-1}^{\ast}\mathcal{I}_{A,-1}=\mathsf P(A)^{-1},
\qquad
\det(\mathcal{I}_{A,-1}^{\ast}\mathcal{I}_{A,-1})=\det\mathsf P(A)^{-1}.\;}
\]
Consequently, the anchored capacity
\(\mathrm{cap}_{-1}(A):=\bigl(\det\mathsf P(A)\bigr)^{\frac{1}{2\kappa}}\)
can be completely rewritten as a standard operator invariant:
\[
\boxed{\;
\mathrm{cap}_{-1}(A)^{-1}
=\bigl\|\wedge^{\kappa}\mathcal{I}_{A,-1}\bigr\|^{1/\kappa}\;}
\]
where \(\wedge^{\kappa}\mathcal{I}_{A,-1}\) denotes the operator induced on the \(\kappa\)-th exterior power, and its norm equals the product of the \(\kappa\) singular values of \(\mathcal{I}_{A,-1}\).
\end{proposition}
\begin{proof}[Proof outline]
By Proposition \ref{prop:xi-anchored-hardy-interpolation-pick-equivalence},
for any \(c\) we have \(\|f_c\|_{H^2}^2=c^\ast\mathsf P(A)^{-1}c\).
Substituting \(f_c=\mathcal{I}_{A,-1}c\) yields
\(\langle c,\mathcal{I}_{A,-1}^{\ast}\mathcal{I}_{A,-1}c\rangle
=c^\ast\mathsf P(A)^{-1}c\), whence \(\mathcal{I}_{A,-1}^{\ast}\mathcal{I}_{A,-1}=\mathsf P(A)^{-1}\).
The exterior power norm identity \(\|\wedge^{\kappa}\mathcal{I}\|^2=\det(\mathcal{I}^{\ast}\mathcal{I})\) is a standard singular value fact for finite-rank operators; substitution yields the stated conclusion. \qedhere
\end{proof}

\begin{corollary}[Covariance Law of the Horizon Stabilizer Subgroup: The Scale Drift of \texorpdfstring{$\mathrm{cap}_{-1}$}{cap-1} Is Sealed by \texorpdfstring{$\phi'(-1)$}{phi'(-1)}]\label{cor:xi-anchored-capacity-stabilizer-covariance}
Let \(\phi\) be a unit disk automorphism with \(\phi(-1)=-1\).
Denote by \(\phi'(-1)>0\) its Julia--Carath\'eodory angular derivative at \(-1\).
Then the pseudo-hyperbolic distance is invariant: \(\rho(\phi(a),\phi(b))=\rho(a,b)\).
Furthermore, the endpoint Poisson weight satisfies
\[
\boxed{\;P_{\phi(a)}(-1)=\phi'(-1)\,P_a(-1)\;}
\]
and consequently, for any finite point set \(A=\{a_j\}_{j=1}^{\kappa}\subset\mathbb{D}\), the covariance law holds:
\[
\boxed{\;
\det\mathsf P(\phi(A))=\phi'(-1)^{\kappa}\det\mathsf P(A),\qquad
\mathrm{cap}_{-1}(\phi(A))=\phi'(-1)^{1/2}\,\mathrm{cap}_{-1}(A).\;}
\]
Therefore, the only cohomological drift of \(-\log\det\mathsf P\) is \(-\kappa\log\phi'(-1)\), while all mutual energy terms (given by \(\rho\)) are strictly invariant.
\end{corollary}
\begin{proof}[Proof outline]
This conclusion is the disk-stabilizer form of Proposition \ref{prop:xi-pick-poisson-det-scaling-law}: on the stabilizer subgroup with \(\phi(-1)=-1\),
\(\rho\) is invariant while \(P_a(-1)\) is covariant with factor \(\phi'(-1)\); then applying the determinant factorization
\(\det\mathsf P=(\prod_j P_{a_j}(-1))(\prod_{j<k}\rho(a_j,a_k)^2)\) (Proposition \ref{prop:xi-pick-poisson-det-factorization}) yields the result. \qedhere
\end{proof}

\begin{corollary}[Extremal Principle for Anchored Capacity: Optimal Upper Bound and Necessary Structure Under Fixed Endpoint Flux]\label{cor:xi-anchored-capacity-extremal-fixed-flux}
Fix \(\kappa\) and let the total endpoint flux be \(p_{\Sigma}:=\sum_{j=1}^{\kappa}p_j\) (where \(p_j=P_{a_j}(-1)\)).
Then
\[
\boxed{\;
\det\mathsf P(A)\ \le\ \Bigl(\frac{p_{\Sigma}}{\kappa}\Bigr)^{\kappa},\qquad
\mathrm{cap}_{-1}(A)\ \le\ \Bigl(\frac{p_{\Sigma}}{\kappa}\Bigr)^{1/2}.\;}
\]
Moreover, if there exists a sequence of configurations for which \(\det\mathsf P(A)\) approaches the above upper bound, then necessarily
(i) the weights become asymptotically uniform: \(p_j\to p_{\Sigma}/\kappa\) (by multiplicity);
(ii) the point pairs achieve asymptotic maximal separation: \(\rho(a_j,a_k)\to 1\) (for all \(j\neq k\)).
In other words, the ``most verifiable'' configuration is jointly forced by \emph{self-energy uniformization} and \emph{mutual energy separation} within the geometric degrees of freedom permitted by the regime.
\end{corollary}
\begin{proof}[Proof outline]
From \(\rho\le 1\) and the determinant factorization, \(\det\mathsf P\le\prod_j p_j\);
then by AM--GM, \(\prod_j p_j\le (p_{\Sigma}/\kappa)^{\kappa}\), yielding the upper bound (the analogous formulation of Corollary \ref{cor:xi-pick-poisson-energy-fixed-flux-extremal}).
If the upper bound is approached, then both the \(\rho\)-term upper bound and the AM--GM equality condition must be simultaneously approached, giving (i) and (ii). \qedhere
\end{proof}

\begin{corollary}[Quantitative stability of anchored capacity extremal: Explicit error bounds for weight uniformization and minimum separation (51)]\label{cor:xi-anchored-capacity-extremal-stability-quant}
Let \(A=\{a_1,\dots,a_\kappa\}\subset\mathbb{D}\) be an anchored configuration, and denote the endpoint fluxes \(p_j:=P_{a_j}(-1)\), \(p_{\Sigma}:=\sum_{j=1}^{\kappa}p_j\).\footnote{Compiled from internal manuscript \cite{Internal2025ResolutionFoldingPhiPi} (2025). Archive timestamp: 2026-02-14 11:15:36 (Asia/Singapore).}
Define the extremal deficit
\[
\boxed{\;
\varepsilon(A):=\kappa\log\!\Bigl(\frac{p_{\Sigma}}{\kappa}\Bigr)-\log\det\mathsf P(A)\ \ge\ 0\;}
\]
Then the following quantitative stability estimates hold:
\begin{enumerate}
  \item Let \(q_j:=p_j/p_{\Sigma}\) (\(\sum_j q_j=1\)) and the uniform distribution \(u_j:=1/\kappa\). Then
  \[
  \boxed{\;
  \sum_{j=1}^{\kappa}|q_j-u_j|
  \ \le\
  \sqrt{\frac{2\,\varepsilon(A)}{\kappa}},\qquad
  \sum_{j=1}^{\kappa}\Bigl|p_j-\frac{p_{\Sigma}}{\kappa}\Bigr|
  \ \le\
  p_{\Sigma}\sqrt{\frac{2\,\varepsilon(A)}{\kappa}}\;}
  \]
  Thus small deficit forces endpoint weights to approximate uniformity in the total variation sense.
  \item Denote the minimum separation \(\rho_{\min}(A):=\min_{j\ne k}\rho(a_j,a_k)\in(0,1]\). Then
  \[
  \boxed{\;
  \rho_{\min}(A)\ \ge\ \exp\!\Bigl(-\frac{\varepsilon(A)}{\kappa(\kappa-1)}\Bigr).\;}
  \]
  In particular, if \(\varepsilon(A)=o(\kappa^2)\), then necessarily \(\rho_{\min}(A)\to 1\), i.e., extremal approximation forces pairwise asymptotic separation.
\end{enumerate}
\end{corollary}

\begin{proof}
By the determinant factorization (Proposition \ref{prop:xi-pick-poisson-det-factorization})
\[
\log\det\mathsf P(A)
=\sum_{j=1}^{\kappa}\log p_j+2\sum_{1\le j<k\le \kappa}\log\rho(a_j,a_k),
 \qquad \rho(a_j,a_k)\in(0,1],
\]
hence
\[
\varepsilon(A)
=\underbrace{\kappa\log\!\Bigl(\frac{p_{\Sigma}}{\kappa}\Bigr)-\sum_{j=1}^{\kappa}\log p_j}_{=:\,\varepsilon_w\ \ge\ 0}
\;+\;
\underbrace{\Bigl(-2\sum_{j<k}\log\rho(a_j,a_k)\Bigr)}_{=:\,\varepsilon_\rho\ \ge\ 0},
\qquad \varepsilon_w,\varepsilon_\rho\le \varepsilon(A).
\]
Let \(q_j=p_j/p_{\Sigma}\), \(u_j=1/\kappa\). Then
\[
\varepsilon_w=\sum_{j=1}^{\kappa}\log\!\Bigl(\frac{p_{\Sigma}/\kappa}{p_j}\Bigr)
=\sum_{j=1}^{\kappa}\log\!\Bigl(\frac{u_j}{q_j}\Bigr)
=\kappa\,D(u\Vert q),
\]
where \(D(u\Vert q)=\sum_j u_j\log(u_j/q_j)\) is the KL divergence. By Pinsker's inequality
\(\|u-q\|_1\le\sqrt{2D(u\Vert q)}\)
we obtain
\[
\sum_{j=1}^{\kappa}|q_j-u_j|
\le
\sqrt{\frac{2\,\varepsilon_w}{\kappa}}
\le
\sqrt{\frac{2\,\varepsilon(A)}{\kappa}},
\]
and multiplying by \(p_{\Sigma}\) yields the second inequality of item 1.
\par
For item 2, since \(\rho(a_j,a_k)\ge\rho_{\min}(A)\) we have
\[
\varepsilon_\rho
=-2\sum_{j<k}\log\rho(a_j,a_k)
\ge
-2\binom{\kappa}{2}\log\rho_{\min}(A)
=\kappa(\kappa-1)\bigl(-\log\rho_{\min}(A)\bigr).
\]
Since \(\varepsilon_\rho\le \varepsilon(A)\), we get
\(-\log\rho_{\min}(A)\le \varepsilon(A)/(\kappa(\kappa-1))\), yielding the stated lower bound. \qed
\end{proof}

\begin{proposition}[Gaussian Mutual Information Formulation of Cross Mutual Energy: The Union Gap Equals the Log-Determinant Mutual Information]\label{prop:xi-cross-energy-gaussian-mutual-information}
Let \(A\sqcup B\) be a disjoint union (counted with multiplicity), and regard \(\mathsf P(A\cup B)\) as a block covariance matrix.
Take a \emph{proper} complex Gaussian vector \(Z_{A\cup B}\sim\mathcal N_{\CC}(0,\mathsf P(A\cup B))\), and let \(Z_A,Z_B\) be its components on the two blocks.
Define the union gap (cross mutual energy)
\[
\mathcal I(A;B):=-\log\det\mathsf P(A\cup B)+\log\det\mathsf P(A)+\log\det\mathsf P(B)\ \ge\ 0.
\]
Then the identity holds
\[
\boxed{\;
\mathcal I(A;B)=I(Z_A;Z_B),\;}
\]
where the right-hand side is the (complex) Gaussian mutual information (in natural logarithms).
Therefore, the cross mutual energy is not only a geometric quantity (Proposition \ref{prop:xi-pick-poisson-union-submultiplicativity-cross-energy}), but also the information-theoretic exact measure of ``two-source residual linkability.''
\end{proposition}
\begin{proof}[Proof outline]
For a proper complex Gaussian vector, the differential entropy formula is \(h(Z)=\log\det(\pi e\,\Sigma)\) (where \(\Sigma\) is the covariance).
Therefore
\(
I(Z_A;Z_B)=h(Z_A)+h(Z_B)-h(Z_{A\cup B})
=\log\det\mathsf P(A)+\log\det\mathsf P(B)-\log\det\mathsf P(A\cup B)
\), yielding the stated identity. \qedhere
\end{proof}

\begin{proposition}[Leading-Order Free Energy of the \texorpdfstring{$\beta$}{beta}-Horizon Ensemble: The \texorpdfstring{$\kappa^2$}{k^2}-Level Is Determined by the Image-Charge Green Energy]\label{prop:xi-beta-ensemble-leading-free-energy}
On the compact feasible domain \(\mathcal K\subset\mathbb H\), consider the \(\beta\)-ensemble
\[
\mathbb P_{\kappa,\beta}(\dd t_1\cdots \dd t_\kappa)
\ \propto\
\exp\!\bigl(-\beta\,S_{\mathrm{gen}}(t_1,\dots,t_\kappa)\bigr)\,\dd t_1\cdots \dd t_\kappa,
\]
where \(S_{\mathrm{gen}}\) is the image-charge Green energy given by Proposition \ref{prop:xi-pick-poisson-image-charge-green-energy} (including self-energy terms).
Denote the partition function by \(Z_{\kappa,\beta}\).
Then, under stability conditions that preclude macroscopic collapse, its leading-order free energy satisfies the formal limit
\[
\boxed{\;
\lim_{\kappa\to\infty}\frac{1}{\kappa^2}\log Z_{\kappa,\beta}
=-\beta\inf_{\mu\in\mathcal P(\mathcal K)}
\iint_{\mathcal K\times\mathcal K}G_{\mathbb H}(t,s)\,\dd\mu(t)\dd\mu(s)\;}
\]
where \(G_{\mathbb H}(t,s)=\log\left|\frac{t-\overline s}{t-s}\right|\).
Moreover, the minimizing equilibrium measure is uniquely characterized by the constant-potential Euler--Lagrange condition (Corollary \ref{cor:xi-pick-poisson-equilibrium-euler-lagrange}).
\end{proposition}
\begin{proof}[Proof outline]
By Proposition \ref{prop:xi-pick-poisson-meanfield-template}, the leading order of \(\kappa^{-2}S_{\mathrm{gen}}\) is given by the double integral energy, while the self-energy terms are only \(O(\kappa)\).
Since \(\log\mathrm{Vol}(\mathcal K^\kappa)=O(\kappa)\), the volume entropy contribution vanishes under \(\kappa^{-2}\) normalization;
hence the leading order is dominated by the minimum of the energy (leading-order form of the Varadhan/Laplace principle). \qedhere
\end{proof}

\begin{conjecture}[Central Limit Theorem for Linear Statistics: Fluctuations of the Anchored Determinantal Point Process Are Controlled by the \texorpdfstring{$H^{1/2}$}{H^{1/2}} Seminorm]\label{conj:xi-kdpp-linear-statistics-clt}
Under the \(\kappa\)-DPP caliber of Proposition \ref{prop:xi-pick-poisson-kdpp}, let \(\kappa\to\infty\) and consider the linear statistic
\(
L_\kappa(\varphi)=\sum_{j=1}^{\kappa}\varphi(t_j)
\)
of a smooth test function \(\varphi\) on a scale where the separation domain and height lower bound are maintained.
Then, after centering and variance normalization, \(L_\kappa(\varphi)\) converges to a standard normal distribution, and the limiting variance is given by the half-order form of the horizon Dirichlet-to-Neumann generator:
\[
\mathrm{Var}(L_\kappa(\varphi))
\ \to\
\frac{1}{2\pi}\int_{\RR}|\xi|\,|\widehat{\varphi_{\partial}}(\xi)|^2\,\dd\xi,
\]
where \(\varphi_{\partial}\) is the corresponding boundary trace (or equivalently, the harmonic projection).
This conjecture aligns the fluctuation structure with the \(|D|\) generator of Proposition \ref{con:discussion-dtn-generator} on the same \(H^{1/2}\) seminorm.
\end{conjecture}

\begin{corollary}[Toeplitz Realization of Negative Spectral Volume: \texorpdfstring{$\det\mathsf P$}{det P} as the Limit of the Finite-Certificate Negative-Subspace Regularized Determinant]\label{cor:xi-toeplitz-negative-regularized-determinant-limit}
Within the stable development interval, let the Toeplitz certificate \(\mathbf{T}_N\) have \(\kappa\) negative eigenvalues \(\{\lambda_j^{-}(\mathbf{T}_N)\}_{j=1}^{\kappa}\).
Define the negative-subspace regularized determinant
\[
\mathrm{Det}_{-}(\mathbf{T}_N):=\prod_{j=1}^{\kappa}\bigl(-\lambda_j^{-}(\mathbf{T}_N)\bigr).
\]
Then in the semiclassical limit \(N\to\infty\) with \(N(1-|a_j|)\to\infty\) (for all defect points),
\[
\boxed{\;
\lim_{N\to\infty}\mathrm{Det}_{-}(\mathbf{T}_N)=\det\mathsf P(A).\;}
\]
\end{corollary}
\begin{proof}[Proof outline]
By Corollary \ref{cor:xi-toeplitz-negative-log-volume-limit},
\(\sum_{j=1}^{\kappa}\log(-\lambda_j^{-}(\mathbf{T}_N))\to \log\det\mathsf P\);
exponentiating both sides yields the stated conclusion. \qedhere
\end{proof}

\begin{corollary}[Explicit Lower Bound for Cross Mutual Energy: Multiplicative Sealing by Minimum Separation]\label{cor:xi-cross-energy-lower-bound-by-min-separation}
For the cross mutual energy \(\mathcal I(A;B)\) of the union \(A\sqcup B\) (Proposition \ref{prop:xi-cross-energy-gaussian-mutual-information}), let
\(
\rho_{\min}(A,B):=\min_{a\in A,b\in B}\rho(a,b)
\).
Then the auditable lower bound holds:
\[
\boxed{\;
\mathcal I(A;B)
=2\sum_{a\in A}\sum_{b\in B}\bigl(-\log\rho(a,b)\bigr)
\ \ge\
2\,|A|\,|B|\cdot\bigl(-\log\rho_{\min}(A,B)\bigr).\;}
\]
\end{corollary}
\begin{proof}[Proof outline]
By the identity of Proposition \ref{prop:xi-pick-poisson-union-submultiplicativity-cross-energy},
\(\mathcal I(A;B)=2\sum_{a\in A,b\in B}(-\log\rho(a,b))\), and using
\(-\log\rho(a,b)\ge -\log\rho_{\min}(A,B)\) yields the result. \qedhere
\end{proof}

\begin{remark}[``Additive--Multiplicative'' Accounting Structure and Two Composition Modes of Anchored Capacity]\label{rem:xi-anchored-capacity-accounting-structures}
Denote \(\Phi(A):=-\log\det\mathsf P(A)=S_{\mathrm{gen}}(A)\) (Definition \ref{def:xi-generalized-defect-entropy-logdet}).
Then for the disjoint union \(A\sqcup B\) there is a strict accounting identity
\[
\Phi(A\cup B)=\Phi(A)+\Phi(B)+\mathcal I(A;B),\qquad \mathcal I(A;B)\ge 0,
\]
where \(\mathcal I(A;B)\) is the cross mutual energy (Proposition \ref{prop:xi-cross-energy-gaussian-mutual-information}).
On the other hand, under tensorized (independent parallel) composition, the exponential quantity \(\det\mathsf P\) is multiplicative while \(\Phi\) is additive.
Thus \(\Phi\) exhibits two structures simultaneously: it is strictly superadditive (with interaction energy term) in the union direction, and strictly additive in the parallel direction.
Combined with the stabilizer-group cohomology drift in Corollary \ref{cor:xi-anchored-capacity-stabilizer-covariance}, \(\Phi\) can be regarded as a canonical accounting potential candidate for the ``horizon data structure.''
\end{remark}
\begin{proof}[Proof Sketch]
By Proposition \ref{prop:xi-pick-poisson-meanfield-template}, the leading order of \(\kappa^{-2}S_{\mathrm{gen}}\) is given by the double integral energy, while the self-energy term is only \(O(\kappa)\).
Since \(\log\mathrm{Vol}(\mathcal K^\kappa)=O(\kappa)\), the volume entropy contribution vanishes under \(\kappa^{-2}\) normalization;
thus the leading order is dominated by the minimum value of energy (leading-order form of Varadhan/Laplace principle). This completes the proof.
\end{proof}

\begin{conjecture}[Central Limit Theorem for Linear Statistics: Fluctuations of Anchored Determinantal Point Process Controlled by \texorpdfstring{$H^{1/2}$}{H^{1/2}} Semi-Norm]\label{conj:xi-kdpp-linear-statistics-clt}
Under the \(\kappa\)-DPP caliber of Proposition \ref{prop:xi-pick-poisson-kdpp}, let \(\kappa\to\infty\) and consider linear statistics of smooth test function \(\varphi\) at scale where separation domain and height lower bound are maintained:
\(
L_\kappa(\varphi)=\sum_{j=1}^{\kappa}\varphi(t_j)
\).
Then after centering and variance normalization, \(L_\kappa(\varphi)\) converges to standard normal distribution, and the limiting variance is given by the half-order form of horizon Dirichlet-to-Neumann generator:
\[
\mathrm{Var}(L_\kappa(\varphi))
\ \to\
\frac{1}{2\pi}\int_{\RR}|\xi|\,|\widehat{\varphi_{\partial}}(\xi)|^2\,\dd\xi,
\]
where \(\varphi_{\partial}\) is the corresponding boundary trace (or equivalently harmonic projection).
This conjecture aligns the fluctuation structure strictly with the \(|D|\) generator of Proposition \ref{con:discussion-dtn-generator} on the same \(H^{1/2}\) semi-norm.
\end{conjecture}

\begin{corollary}[Toeplitz Realization of Negative Spectrum Volume: \texorpdfstring{$\det\mathsf P$}{det P} as Limit of Regularized Determinant of Finite Certificate Negative Subspace]\label{cor:xi-toeplitz-negative-regularized-determinant-limit}
Within stable development interval, let Toeplitz certificate \(\mathbf{T}_N\) have \(\kappa\) negative eigenvalues \(\{\lambda_j^{-}(\mathbf{T}_N)\}_{j=1}^{\kappa}\).
Define negative subspace regularized determinant
\[
\mathrm{Det}_{-}(\mathbf{T}_N):=\prod_{j=1}^{\kappa}\bigl(-\lambda_j^{-}(\mathbf{T}_N)\bigr).
\]
Then in the semiclassical limit \(N\to\infty\) with \(N(1-|a_j|)\to\infty\) (for all defect points),
\[
\boxed{\;
\lim_{N\to\infty}\mathrm{Det}_{-}(\mathbf{T}_N)=\det\mathsf P(A).\;}
\]
\end{corollary}
\begin{proof}[Proof Sketch]
By Corollary \ref{cor:xi-toeplitz-negative-log-volume-limit},
\(\sum_{j=1}^{\kappa}\log(-\lambda_j^{-}(\mathbf{T}_N))\to \log\det\mathsf P\);
taking exponentials on both sides yields the stated conclusion. This completes the proof.
\end{proof}

\begin{corollary}[\texorpdfstring{$\kappa^2$}{k^2}-Order Determinant Collapse for Macroscopic Splitting: Exponential Compression Law from Cross Mutual Energy Lower Bound (52)]\label{cor:xi-cross-energy-lower-bound-by-min-separation}
For the cross mutual energy \(\mathcal I(A;B)\) of union \(A\sqcup B\) (Proposition \ref{prop:xi-cross-energy-gaussian-mutual-information}), denote
\(
\rho_{\min}(A,B):=\min_{a\in A,b\in B}\rho(a,b)
\).
Then the auditable lower bound holds
\[
\boxed{\;
\mathcal I(A;B)
=2\sum_{a\in A}\sum_{b\in B}\bigl(-\log\rho(a,b)\bigr)
\ \ge\
2\,|A|\,|B|\cdot\bigl(-\log\rho_{\min}(A,B)\bigr).\;}
\]
Thus from the definition of \(\mathcal I(A;B)\), the strict compression law follows
\[
\boxed{\;
\det\mathsf P(A\cup B)
=\det\mathsf P(A)\det\mathsf P(B)\,e^{-\mathcal I(A;B)}
\ \le\
\det\mathsf P(A)\det\mathsf P(B)\,\rho_{\min}(A,B)^{2|A||B|}\; }.
\]
In particular, if there exists constant \(\rho_0<1\) such that \(\rho_{\min}(A,B)\le \rho_0\), and \(|A|\) and \(|B|\) grow at same order, then the relative collapse ratio satisfies
\(
\det\mathsf P(A\cup B)\big/(\det\mathsf P(A)\det\mathsf P(B))
\le \exp(-c\,|A||B|)
\),
where \(c=-2\log\rho_0>0\), thus exponential collapse at \(\exp(-\Theta(\kappa^2))\) level appears in macroscopic splitting regime.
\end{corollary}
\begin{proof}[Proof Sketch]
From the identity \(\mathcal I(A;B)=2\sum_{a\in A,b\in B}(-\log\rho(a,b))\) in Proposition \ref{prop:xi-pick-poisson-union-submultiplicativity-cross-energy}, using
\(-\log\rho(a,b)\ge -\log\rho_{\min}(A,B)\) yields the first inequality.
Rewriting \(\mathcal I(A;B)=-\log\det\mathsf P(A\cup B)+\log\det\mathsf P(A)+\log\det\mathsf P(B)\) as
\(
\det\mathsf P(A\cup B)=\det\mathsf P(A)\det\mathsf P(B)e^{-\mathcal I(A;B)}
\),
and substituting \(\mathcal I(A;B)\ge 2|A||B|(-\log\rho_{\min})\) yields the compression law. This completes the proof.
\end{proof}

\begin{corollary}[Determinant Exponential Collapse for Multi-Block Splitting: Full Partition Version of \texorpdfstring{$\kappa^2$}{k^2}-Order Compression]\label{cor:xi-pick-poisson-det-multipartition-kappa2-collapse}
Let finite anchored set \(A\) decompose into disjoint union \(A=\biguplus_{r=1}^{R}A_r\), and denote \(\kappa_r:=|A_r|\).
Under the determinant factorization regime of Proposition \ref{prop:xi-pick-poisson-det-factorization}, the exact factorization always holds
\[
\boxed{\;
\det\mathsf P(A)
=\Bigl(\prod_{r=1}^{R}\det\mathsf P(A_r)\Bigr)\,
\prod_{1\le r<s\le R}\ \prod_{a\in A_r,\ b\in A_s}\rho(a,b)^{2}\; }.
\]
Thus the unified upper bound follows: letting
\(
\rho_{\min}(A_r,A_s):=\min_{a\in A_r,\ b\in A_s}\rho(a,b)
\),
\[
\boxed{\;
\det\mathsf P(A)\ \le\
\Bigl(\prod_{r=1}^{R}\det\mathsf P(A_r)\Bigr)\,
\prod_{1\le r<s\le R}\rho_{\min}(A_r,A_s)^{2\kappa_r\kappa_s}\; }.
\]
In particular, if there exists constant \(\rho_0<1\) such that \(\rho_{\min}(A_r,A_s)\le\rho_0\) for all \(r<s\), then
\[
\det\mathsf P(A)\ \le\ \Bigl(\prod_{r=1}^{R}\det\mathsf P(A_r)\Bigr)\,
\exp\!\Bigl(-2(-\log\rho_0)\sum_{1\le r<s\le R}\kappa_r\kappa_s\Bigr),
\]
exhibiting universal \(\kappa^2\)-order exponential collapse law.
\end{corollary}
\begin{proof}
Decompose \(\prod_{a<b}\rho(a,b)^2\) into "intra-block pairs" and "cross-block pairs": intra-block pairs correspond exactly to \(\prod_r\det\mathsf P(A_r)/\prod_{a\in A}P_a(-1)\), cross-block pairs remain explicit products, yielding the first formula.
The second formula follows immediately by applying \(\rho(a,b)\le \rho_{\min}(A_r,A_s)\) term by term to cross-block pairs; the last formula follows by taking logarithm and rearranging. This completes the proof.
\end{proof}

\begin{remark}[``Additive--Multiplicative'' Accounting Structure of Anchored Capacity and Two Composition Modes]\label{rem:xi-anchored-capacity-accounting-structures}
Denote \(\Phi(A):=-\log\det\mathsf P(A)=S_{\mathrm{gen}}(A)\) (Definition \ref{def:xi-generalized-defect-entropy-logdet}).
Then for the union \(A\sqcup B\), the strict accounting identity holds:
\[
\Phi(A\cup B)=\Phi(A)+\Phi(B)+\mathcal I(A;B),\qquad \mathcal I(A;B)\ge 0,
\]
where \(\mathcal I(A;B)\) is the cross mutual energy (Proposition \ref{prop:xi-cross-energy-gaussian-mutual-information}).
On the other hand, under tensorization (independent parallel composition), the exponential-type quantity \(\det\mathsf P\) is multiplicative, while \(\Phi\) is additive.
Therefore, \(\Phi\) simultaneously exhibits two structures: strict superadditivity in the union direction (with interaction energy terms), and strict additivity in the parallel direction.
Combined with the stabilizer subgroup cohomological drift of Corollary \ref{cor:xi-anchored-capacity-stabilizer-covariance}, \(\Phi\) can be regarded as a candidate canonical bookkeeping potential for the ``horizon data structure.''
\end{remark}

\begin{conjecture}[Capacity Collapse Criterion for Continuous Defects: Continuous Hair Forces Discrete Anchored Capacity to Zero]\label{conj:xi-continuous-hair-capacity-collapse}
Let a general defect be induced by the Herglotz/Clark measure decomposition, and suppose it contains a nontrivial singular inner factor (continuous hair).
Then for any discrete approximation family \(A_\kappa\subset\mathbb D\) (\(\kappa\to\infty\)), if its visible boundary density locally approximates this continuous defect, then necessarily
\[
\mathrm{cap}_{-1}(A_\kappa)\ \longrightarrow\ 0.
\]
This conjecture echoes Corollary \ref{cor:xi-discrete-vs-continuous-hair-by-negative-inertia}: ``unbounded negative inertia \(\Leftrightarrow\) continuous hair/infinite point-wise degrees of freedom'': a continuous defect cannot be discretized with positive capacity under a fixed anchored caliber, and its volume-type verifiable surplus must be reclaimed to zero in the limit.
\end{conjecture}

\endinput




\endinput


\paragraph{Basepoint scanning tomography: shifted Cayley's Cauchy--Poisson profile determines full off-axis zero parameters}
In fixed charts, off-critical-slice offset undergoes quadratic redshift compression in high-height intervals (Theorem \ref{thm:xi-offcritical-quadratic-radial-compression}), while the comoving chart of shifted Cayley can remove this compression (\S\ref{subsubsec:xi-comoving-defect-delta-bound}).
The following result elevates this "comoving scanning" to strict inversion: a one-dimensional scanning profile carries all discrete degrees of freedom of two-dimensional zero parameters.

\begin{definition}[Basepoint scanning profile]\label{def:xi-basepoint-scan-profile}
Let \(\rho_k=\beta_k+\mathrm{i}\gamma_k\) (\(1\le k\le\kappa\)) be a group of off-critical-line zeros (counted with multiplicity), and denote \(\delta_k:=\tfrac12-\beta_k>0\).
For \(x_0\in\RR\), let its shifted disk image \(w_{\rho_k,x_0}\) be as defined in \eqref{eq:xi-comoving-disk-image}, and let
\[
h_{x_0}(\rho_k):=1-|w_{\rho_k,x_0}|^2
\qquad(\text{comoving depth}).
\]
Define the basepoint scanning profile as
\[
\boxed{\;
H(x_0):=\sum_{k=1}^{\kappa}h_{x_0}(\rho_k)\; }.
\]
\end{definition}

\begin{proposition}[Closed form of comoving depth]\label{prop:xi-basepoint-depth-closed-form}
Under the setting of Definition \ref{def:xi-basepoint-scan-profile}, for each \(k\) there holds the strict identity
\[
\boxed{\;
h_{x_0}(\rho_k)=\frac{4\delta_k}{(\gamma_k-x_0)^2+(1+\delta_k)^2}\; }.
\]
\end{proposition}

\begin{proof}
By \eqref{eq:xi-comoving-disk-image}, \(w_{\rho_k,x_0}=\mathcal{C}((\gamma_k-x_0)+\mathrm{i}\delta_k)\).
Applying the modulus gap formula for Cayley compactification \(1-|\,\mathcal{C}(\gamma+\mathrm{i}\delta)\,|^2=\frac{4\delta}{\gamma^2+(1+\delta)^2}\) (closed form part of Theorem \ref{thm:xi-offcritical-quadratic-radial-compression}) to \(\gamma\mapsto \gamma_k-x_0\) yields the result. Proven.
\end{proof}

\begin{theorem}[Basepoint scanning inversion uniqueness: \texorpdfstring{$H$}{H} determines \texorpdfstring{$\{(\gamma_k,\delta_k)\}$}{(gamma,delta)}]\label{thm:xi-basepoint-scan-inversion-uniqueness}
Under finite defect setting, \(H\) is a real analytic function of \(x_0\), and has rational meromorphic continuation on the complex plane:
its pole set (with multiplicity) is exactly
\[
x_0=\gamma_k\pm \mathrm{i}(1+\delta_k)\qquad(1\le k\le\kappa).
\]
Therefore, given the values of \(H\) on any nontrivial real interval, one can uniquely determine the multiset \(\{(\gamma_k,\delta_k)\}\) (counted with multiplicity).
In other words, off-axis zeros appear in the scanning variable \(x_0\) as strict Cauchy--Poisson peaks: peak center locates \(\gamma_k\), peak half-width carries \((1+\delta_k)\); the comoving choice \(x_0=\gamma_k\) decouples \(\delta_k\) from quadratic redshift compression and reveals it at linear scale (see \eqref{eq:xi-comoving-pure-delta}).
\end{theorem}

\begin{proof}
By Proposition \ref{prop:xi-basepoint-depth-closed-form},
\(
H(x_0)=\sum_k \frac{4\delta_k}{(\gamma_k-x_0)^2+(1+\delta_k)^2}
\).
Each term is a constant multiple of upper half-plane Poisson kernel, hence is the boundary imaginary part of some rational Herglotz function; viewing \(x_0\) as complex variable gives its meromorphic continuation, with denominator vanishing yielding pole locations \(x_0=\gamma_k\pm \mathrm{i}(1+\delta_k)\).
If two parameter groups yield the same \(H\) on a real interval, then the difference function vanishes on that interval and is the boundary value of a meromorphic function; by analytic continuation uniqueness, its meromorphic extension is identically zero, thus pole sets (with multiplicity) coincide, whence \(\{(\gamma_k,\delta_k)\}\) coincide. Proven.
\end{proof}

\paragraph{Equivariant obstruction triangular alignment: unification of \texorpdfstring{$\deg(B)$}{deg(B)}, Maslov index, and Pontryagin negative squares}
\begin{theorem}[Triple unification of point obstruction]\label{thm:xi-tate-maslov-pontryagin-unification}
Under the finite defect and no singular inner factor framework of this section, the point equivariant obstruction is the same integer \(\kappa\), and
\[
\boxed{\;
\kappa=\deg(B_{F_{\zeta}})
=\dim K_{B_{F_{\zeta}}}
=\dim\mathcal{H}_{\mathrm{pp}}(A_{B_{F_{\zeta}}})
=\nu_-(\mathbf{T}_N)\ \ (N\gg 1)\;}
\]
where \(\dim K_{B_{F_{\zeta}}}=\mathfrak d_{\mathrm{Bl}}(F_{\zeta})\) is given by Definition \ref{def:blaschke-defect-dimension},
\(\dim\mathcal{H}_{\mathrm{pp}}(A_{B_{F_{\zeta}}})=\deg(B_{F_{\zeta}})\) is given by Theorem \ref{thm:blaschke-point-spectrum-correspondence},
and the stability \(\nu_-(\mathbf{T}_N)=\kappa\) is given by Theorem \ref{thm:xi-toeplitz-negative-squares-equals-blaschke-defect}.
In the \(\ZZ/2\)-equivariant uniformization semantics, this \(\kappa\) is exactly the point-like (Maslov/Toeplitz/Fredholm) integer part of the Tate obstruction (Proposition \ref{prop:selfdual-maslov-singular-dichotomy}), while the continuous part is carried by the singular inner factor \(S_{F_{\zeta}}\); under the \(S_{F_{\zeta}}\equiv 1\) regular framework of this section, the continuous part vanishes.
\end{theorem}

\endinput
